\documentclass[reqno]{amsart}
%\documentclass[letter,12pt]{report}
%\usepackage{amsaddr}
%\documentclass[nofootinbib,report,12pt]{revtex4-2}

\usepackage[margin=3cm]{geometry} %set the margins

\usepackage{bm}
\usepackage{hyperref} %[hidelinks] hides the fluorescent boxes and makes all colors black
\hypersetup{
    colorlinks = true,
    linkcolor={gray},
    citecolor={blue!50!black},
    urlcolor={blue!50!black}
}
\usepackage[utf8]{inputenc}
\usepackage{cite}
%\usepackage[sort&compress]{natbib}

\usepackage{amsmath,accents}
%\numberwithin{equation}{section} %numbers the equations according to each section.

\usepackage{siunitx}
\usepackage{mathtools}
\usepackage{amssymb}
\usepackage[normalem]{ulem}
\usepackage{enumitem}
%\graphicspath{{figures/}}

\usepackage{tcolorbox} %allows one to put colored boxes around text and equations (cannot be used with xcolor package}

% For tables
\usepackage{booktabs,siunitx,array,threeparttable}
\sisetup{group-minimum-digits=4}

\def\be{\begin{equation}}
\def\ee{\end{equation}}
\def\e#1{\label{#1}\end{equation}}
\def\bea{\begin{eqnarray}}
\def\eea{\end{eqnarray}}
\def\ea#1{\label{#1}\end{eqnarray}}
\def\bse{\begin{subequations}}
\def\ese{\end{subequations}}
\def\nn{\nonumber}

\def\r#1{(\ref{#1})}
\def\bem#1{\begin{mathletters}\label{#1}}
\def\eml{\end{mathletters}}
\def\df{\Delta\varphi}
\def\agt{\stackrel{>}{_\sim}}
\def\unc#1{\delta \tilde{#1}}
\def\uncwrt#1#2{\delta_{\mathrm{#2}} \tilde{#1}}
\def\uncwrtpost#1#2#3{\delta_{\mathrm{#2}} \tilde{#1}_{(#3)}}
\def\est#1{\tilde{#1}}

\newcommand{\ket}[1]{\left| #1 \right>}
\newcommand{\bra}[1]{\left< #1 \right|}
\newcommand{\hham}{\hat{H}}
\newcommand{\U}{\hat{U}}
\newcommand{\iu}{{i\mkern1mu}}
\def\braket#1#2{{\langle#1|#2\rangle}}
\def\mean#1{{\langle#1\rangle}}
\newcommand{\myvec}[1]{\accentset{\rightharpoonup}{#1}}
\def\4#1{{\boldsymbol{#1}}}
\def\8#1{\mathbb{E}\left[#1\right]}
\def\func#1#2{\mathrm{#1}\left[#2\right]}
\def\funcwrt#1#2#3{\mathrm{#1}_{\mathrm{#3}}\left[#2\right]}

\DeclareMathSymbol{\mlq}{\mathord}{operators}{``}
\DeclareMathSymbol{\mrq}{\mathord}{operators}{`'}

%\renewcommand\thefootnote{\textcolor{gray}{\arabic{footnote}}} %Not needed if hypersetup is used

\usepackage{titlesec} %titlesec package
\titleformat{\subsection}{\centering\normalfont\bf}{\thesubsection.\quad}{0em}{}{}
\renewcommand{\thesubsection}{\arabic{subsection}}
%\renewcommand{\theequation}{\arabic{section}.\arabic{equation}}

%Allows me to make lists with custom enumeration and also to link a label to that enumeration
\makeatletter
\newcommand{\myitem}[1]{%
\item[#1]\protected@edef\@currentlabel{#1}%
}
\makeatother

\begin{document}

\title{Quantum squeezing cannot beat the standard quantum limit}
%\title{Squeezed states cannot beat the standard quantum limit or the single particle limit}

\maketitle

\begin{centering}
\author{Liam P. McGuinness}\\
\address{Laser Physics Centre, Research School of Physics, Australian National University, Acton,}\\
\address{Australian Capital Territory 2601, Australia}\\
\email{Email: \href{mailto:liam@grtoet.com}{liam@grtoet.com}}\\
\end{centering}

%\author{Liam P. McGuinness}
%\email[Email: ]{liam@grtoet.com}
%\affiliation{Laser Physics Centre, Research School of Physics, Australian National University, Acton, Australian Capital Territory 2601, Australia}


\begin{abstract}
Quantum entanglement between particles is expected to allow one to perform tasks that would otherwise be impossible \cite{Bell1964,Feynman1982,Deutsch1985, Shor1994, Giovannetti2001}. In quantum sensing and metrology, entanglement is often claimed to enable a precision that cannot be attained with the same number of particles and time, forgoing entanglement \cite{Helstrom1969, Holevo1973a, Giovannetti2004, Giovannetti2006, DemkowiczDobrzanski2012, Zwierz2012, Pezze2018}. Two distinct approaches exist: creation of entangled states that either \emph{i)} respond quicker to the signal, or \emph{ii)} are associated with lower noise and uncertainty. The second class of states are generally called squeezed states. Here we show that if our definition of success is -- a precision that is impossible to achieve using the same resources but without entanglement -- then squeezed states cannot succeed. In doing so we show that a single non-separable squeezed state provides fundamentally no better precision, per unit time, than a single particle.\\
\end{abstract}

\section*{Prelude}
I have asked for and received a lot of feedback on this work from experts in the field\footnote{If you have any feedback, please contact me.}. Here I try to distil the arguments presented in the main text as clearly as possible.\\

If one wants to compare the precision of two measurement devices, a good method is to break each device down into its individual components and analyse the precision each component provides. If each component of the first device provides no better precision than each component of the second, it is possible to conclude that the first device cannot outperform the second. For rigour, two additional points are needed:
\begin{enumerate}
\item A check that the second device has more components (as more components improve precision).
\item An assumption that the components are independent. 
\end{enumerate}

This paper performs such an analysis. The mathematical content is simply to analyse the (optimum) precision per unit time, that quantum mechanics predicts an individual component can achieve\footnote{To be precise, I analyse the maximum information each individual component provides per unit time, for a single measurement.}. One proviso is that this analysis only applies when each component is measured once. The physics is to associate an individual component of each device with a single non-separable state vector. As entangled particles form non-separable states, this analysis of the measurement precision accounts for quantum correlations. However, with assumption (2) above, these are the only form of correlations present. 

Using basic quantum mechanics, it is straightforward to bound the amount of information per unit time, a single state vector can provide on a signal $\theta$. The amount of information (and the measurement precision) just depends on the state vector response to $\theta$, i.e. $\frac{\mathrm{d}}{\mathrm{d}\theta}$. This is a well-known and standard result (see e.g. Refs.\,[\textcolor{blue!50!black}{91}--\textcolor{blue!50!black}{93}]). In terms of the mathematics and physics there is nothing even mildly controversial about what I am saying here. Note, we don't even need to worry about the length of the state vector -- entangled or unentangled -- any pure state has unit length.

I then make the observation that in the quantum squeezing community, people claim to improve the precision of their measurement device (beyond a fundamental limit) without improving the state response to the signal. In fact, in a squeezed device, the response to the signal of a single non-separable squeezed state is often described as being the same as that of a single particle. In case you think that I am misrepresenting the position of the squeezing community, please look through the papers for yourself. The review papers included as Refs.\,[\textcolor{blue!50!black}{12}--\textcolor{blue!50!black}{15}] (see also the references within [\textcolor{blue!50!black}{14}]) and Refs.\,[\textcolor{blue!50!black}{33},\,\textcolor{blue!50!black}{35}--\textcolor{blue!50!black}{85},\,\textcolor{blue!50!black}{96}--\textcolor{blue!50!black}{103}] are a good starting point. There is a mathematical and logical contradiction between these two statements -- they cannot both be true. We cannot simultaneously believe that: \textit{a}) the optimum precision of any non-separable state vector can be characterised by the state response to a parameter, and \textit{b}) squeezed states can improve the optimum precision without improving the state response. It turns out that the results of the squeezing community are in error.

To prove that a device constructed from an ensemble of $N$ squeezed particles can not provide fundamentally more information per unit time than possible from a device made from the same number of independent particles (this is known as the standard quantum limit), I note:
\begin{enumerate}
\item For an ensemble containing $N > 1$ particles, if the ensemble is entangled to any degree, it must contain less indivisible state vectors than an unentangled ensemble.
\end{enumerate}
That's it. Much of the main text discusses experiments with spins, but the analysis is general and also applies to measuring devices composed of photons, i.e. optical interferometers such as LIGO.\\

\noindent Just to reiterate the logical arguments made here.
\begin{enumerate}
\item Assuming the postulates of quantum mechanics are correct, the minimum uncertainty per unit time, that any single non-separable state provides on an unknown signal is bounded by the state response to the signal. This bound takes into account the minimum noise when measuring a single state vector.
\item In the literature, a squeezed state is commonly described as having the same signal response as an unentangled single particle state. Instead of an increased response speed, the noise when measuring this state is claimed to be less than that achievable when measuring a single particle. Thus the claim is that squeezed states can achieve an enhanced precision. I agree that this argument sounds reasonable. It is persuasive because it appeals to our intuitive concept of signal-to-noise determining the measurement precision. I refer you to preceding point. Based on standard quantum mechanics, this is impossible.
\item For fixed number of particles, an ensemble containing entanglement has less non-separable state vectors than an unentangled ensemble. As squeezed states are entangled, measuring devices constructed using squeezed ensembles cannot beat the standard quantum limit.\\
\end{enumerate}

If I could pinpoint why people have such strong resistance to these results, it seems to be because there is a vast body of theoretical and experimental work claiming the contrary. If you are at that point, then I ask you first to go through the arguments that I am making and look for an error. If none is forthcoming, then I urge you to consider reassessing the analysis presented by the rest of the scientific community (see the FAQ's section and Ref.\,[\textcolor{blue!50!black}{31}] for a summary of my efforts to date).

I have lost a lot of confidence in the peer-review system (see Refs.\,[\textcolor{blue!50!black}{114}--\textcolor{blue!50!black}{116}]), but that is not to say that I do not want feedback. Of course I would like other people to check my work for rigour. Not only that, I also want to demonstrate to non-experts that this has been done. I do not think that the peer-review system performs those two tasks very well. For those two goals, a much more effective strategy is to financially incentivise people to perform a thorough peer-review, i.e. offer a financial reward to anyone who finds an error. Therefore I offer a prize of US\$10,000 to the first person that finds an error that invalidates the conclusions of this work. Please send to me via email\footnote{Note, in Ref.\,[\textcolor{blue!50!black}{31}] I have also offered a prize of US\$10,000 for the first experimental demonstration of a precision beyond the standard quantum limit, and an additional prize of US\$10,000 for the first demonstration using $N$-partite entanglement that surpasses the single particle limit (in defining the resources, the measurement time and number of particles should at the least include all unitary evolution). I have emailed the authors of several papers claiming to surpass the SQL and offered the prize if they send me the data substantiating their claims (e.g. Refs.\,[\textcolor{blue!50!black}{84, 85, 101}]). Despite this, the prize has remained unclaimed after nearly two years. }.% In addition, I would like to hear from anyone who agrees with the conclusions of this paper. In fact, I will pay US\$100 each, to the first ten people who write to me stating that they agree with these conclusions, and give at least one physical justification as to why.}.

\newpage

\section*{Introduction}
%We first give a non-rigorous discussion of measurement precision in the context of spin-1/2 particles and outline arguments towards the proof. % a layperson explanation of the proof.\\
Consider a sensing device composed of $N$ spin-$1/2$ particles\footnote{Note, these results apply to photon interferometry and are discussed in this context in Appendix\,\ref{App:LIGO}.}, which is used to measure some unknown signal $\theta$. %Typical scenarios include measuring the amplitude or direction of a magnetic field, frequency estimation or phase estimation. 
Taking just one spin and using the spin direction $\myvec{S}$ of this particle as a meter, then by error propagation any estimate $\est{\theta}$ of the value of $\theta$ has an uncertainty\footnote{More fully, we can include the possibility of a finite bias $\epsilon$ and average over $p(\theta = \Theta)$, the \emph{a priori} probability distribution of $\theta$, to obtain the expected uncertainty: $\langle \Delta \est{\theta}\rangle \geq \sqrt{\int^{\Theta} p(\theta = \Theta)\left( \left| \frac{\partial \myvec{S}}{\partial \theta} \right|^{-2} \left(\Delta \myvec{S}\right)^2+\epsilon^2\right) d\theta}$.}:
\be
\Delta \est{\theta} \geq \left| \frac{\partial \myvec{S}}{\partial \theta} \right|^{-1} \Delta \myvec{S}, % + \mathrm{Bias},
\label{eq:SNR}
\ee
where $\Delta \myvec{S}$ is the uncertainty in determining the spin direction. The derivative term is often called the measurement signal, and $\Delta \myvec{S}$ %$(\Delta \myvec{S})^{-1}$ 
the measurement noise\footnote{Incorrectly so, it is the measurement uncertainty not noise.}, so that the `signal-to-noise ratio' determines the measurement precision. With $N$ identical and independent spins, in general one cannot do better than
\be
\Delta \est{\theta} \geq \Delta \est{\theta}_1 / \sqrt{N}, \label{eq:SQL}
\ee
where $\Delta \est{\theta}_1$ bounds the uncertainty using a single particle. Note, for Eq.\,\eqref{eq:SQL} to hold, $\Delta \est{\theta}_1$ must represent the best possible uncertainty that can be attained with a single particle, otherwise one could beat the limit simply by improving $\Delta \est{\theta}_1$. Furthermore, in general $\Delta \est{\theta}_1$ is a function of time since the spin response is given by unitary evolution and with more time one can increase the spin response or perform more measurements to improve the uncertainty. Under this definition, Eq.\,\eqref{eq:SQL} is called the standard quantum limit (SQL) and it sets an uncertainty bound per unit time that is impossible to surpass with a given number of identical and independent spins.

There are two approaches to overcoming the SQL using entanglement, either make the spin response greater, or reduce the uncertainty in measuring the spin direction \cite{Gross2012, Ma2011, Ludlow2015, Pezze2018}. The first makes use of entangled NOON, CAT or GHZ states so that, in theory: $\frac{\partial \myvec{S}_N(t)}{\partial \theta} = N \frac{\partial \myvec{S}_1(t)}{\partial \theta}$, whilst $\Delta \myvec{S}_N = \Delta \myvec{S}_1$, where the subscript denotes the number of spins in the sensor \cite{Pezze2018, Gross2012, Ludlow2015, Bollinger1996, Dowling1998, Childs2000, Campos2003, Kok2004, Leibfried2004, Nagata2007, Resch2007, Berry2009, Dorner2009, Escher2011, Giovannetti2011, Daryanoosh2018, Shaniv2018, Xie2021}. In effect a single spin with greater magnetic moment is created such that the response of the entangled device is $N$-fold greater than that of a single spin. It is worth noting here, that $\Delta \myvec{S}$ is the result of sampling from an unknown probability distribution and since a measurement of both a single entangled state and a single spin provide only one sample, they have the same measurement uncertainty. We will not discuss this approach except to note that when the resources required to generate entanglement are fully accounted for, no improvement over the SQL or $\Delta \est{\theta}_1$ has been demonstrated \cite{McGuinness2021}.

The second approach uses entangled squeezed states so that\footnote{Early squeezing proposals did not require entanglement to overcome the SQL, resulting in confusion as to whether entanglement is required and the definition of a squeezed state (see Appendix\,\ref{App:History}). This definition is still not uniformly standardized see \cite{Soerensen2001,Meyer2001,Maccone2020}, and conflicting definitions of the SQL also remain. We explicitly define the SQL as whatever one can achieve without entanglement, thus it is a moot point whether entanglement is required to surpass it.}: $\frac{\partial \myvec{S}_N(t)}{\partial \theta} = \frac{\partial \myvec{S}_1(t)}{\partial \theta}$, whereas $\Delta \myvec{S}_N = \Delta \myvec{S}_1/N$. I.e. the spin response of the entangled device remains the same as that of a single spin but the measurement noise reduces by a factor of $N$ \cite{Caves1981, Ma2011, Pezze2018, Gross2012, Ludlow2015, Walls1983, Wodkiewicz1985, Wu1986, Slusher1987, Xiao1987, Polzik1992, Wineland1992, Kitagawa1993, Wineland1994, Sanders1995, Kuzmich1998, Soerensen1998, Brif1999, Kuzmich2000, Meyer2001, Esteve2008, Appel2009, Eberle2010, Gross2010, Koschorreck2010, Koschorreck2010a, Leroux2010, Leroux2010a, LouchetChauvet2010, SchleierSmith2010, Wasilewski2010, Luecke2011, Sewell2012, Aasi2013, Taylor2013, Bohnet2014, Muessel2014, Strobel2014, Kruse2016, Polzik2016, Cox2016, Davis2016, Bohnet2016, Hosten2016, Linnemann2016, Macri2016, Tse2019, Braverman2019, Schulte2020, Malia2020, Bao2020, PedrozoPenafiel2020, Casacio2021, Gilmore2021, Greve2022, Malia2022}. We can already see that there is a conflict between our explanation of the origin of the measurement uncertainty and what is observed with squeezed states. Shouldn't the uncertainty in estimating the direction of a single squeezed state be the same as for a single spin? Where does this reduced uncertainty come from? Either our layperson description is wrong or squeezed states do not perform as described. We now show that, if entanglement provides any benefit over the SQL, then it must come about from increasing the sensor response to the signal and not through reduced noise\footnote{My apologies for changing between `uncertainty' and `noise', they do not mean the same thing. I would like to use the correct terminology - uncertainty, but the papers in this field refer to noise and I am quoting their claims. This conflation of terms is the crux of the issue, I encourage the reader to take careful note of this point. Returning to the start of this paragraph, $\Delta \myvec{S}_N = \Delta \myvec{S}_1/N$ does not mean the measurement noise reduces by a factor of $N$, rather the uncertainty in estimating the spin direction of $\myvec{S}_N$, is $N$ times lower than estimating the direction of a single particle $\myvec{S}_1$.}.\\

\section*{Noise independent quantum precision bound}% Quantum precision bound only depends on detector response to signal.}
For rigour we make two adjustments to our language. First, rather than talk about spin direction, we address the underlying mathematical object -- the state vector. To show that the uncertainty bound obtained from an entangled ensemble containing squeezed states is worse than that of an unentangled ensemble, we use a counting argument based on the number of indivisible state vectors (i.e. a state vector that cannot be further factored into non-separable states) in the ensemble. Henceforth, a state vector, $\ket{\psi}$ always refers to this basic unit we are dealing with -- an indivisible state. A single state vector is never used to describe the quantum state of a separable ensemble, instead we keep note of the number of copies of each state. Secondly, for technical reasons and conciseness we avoid quantitatively defining uncertainty, we define instead the (Fisher) information on $\theta$ denoted $\func{I}{\theta,t}$, provided by a given state, or more precisely, measurement of that state \cite{Braunstein1994}:
\be
\func{I}{\theta,t} \equiv \int^{X} \mathrm{d}X \frac{1}{\func{Pr}{X|\theta,t}}\left(\frac{\partial \func{Pr}{X|\theta,t}}{\partial \theta}\right)^2, \quad \func{I}{\theta} \equiv \sum_{i=1}^R \frac{1}{\func{Pr}{X_i|\theta,t}}\left(\frac{\partial \func{Pr}{X_i|\theta,t}}{\partial \theta}\right)^2
\label{eq:QFI}
\ee
where $\func{Pr}{X|\theta,t}$ is the conditional probability to obtain the measurement result $X$ in time $t$ given $\theta$ and the (LHS) RHS assumes the measurement outcomes are (continuous) discrete with $R$ possibilities. We note that for any measurement, the estimation uncertainty $\Delta \est{\theta}$ is a non-increasing function of $\func{I}{\theta,t}$. This observation provides the necessary tools to compare uncertainties, to say one is greater than another, and is sufficient for our purposes.\\

\emph{Key requirement:} The central claim in squeezing enhanced metrology is that squeezed states have an improved intrinsic noise (uncertainty) compared to a single spin. It is clear that the response to $\theta$ of the squeezed state, quantified by $\left|\frac{d\ket{\psi(\theta,t)}}{d \theta}\right|$, is not greater than that of a single spin. One reason this must be the case is that if the noise can be reduced by a factor of $N$, then any further improvement in the sensor response would violate the Heisenberg limit.\\% Our proof directly addresses this claim, we show an information bound that just depends on the state response to $\theta$, therefore, if the state response is not increased, the uncertainty cannot be reduced beyond the SQL.

Finally, we only establish a bound on the information provided by a single measurement of the ensemble. While this allows a direct like-for-like comparison between measurement devices, it is important to note that this is not how measurement sensitivity is usually reported. Often, when reporting squeezing enhancements, comparison is made to a time-averaged limit with $\mathrm{Hz}^{-1/2}$ uncertainty improvement. Further assumptions for the proof are provided in Appendix\,\ref{App:Assumptions}.\\
%\begin{itemize}
%\item A single measurement of each particle in the ensemble.
%\item Where the path parameterized by $\theta$ is smooth. I.e. the probability distribution of $\theta$ is smooth.
%\item $\left(\Delta \est{\theta}\right)^2$ denotes the expected mean squared error of the estimator, thus $\Delta \est{\theta}$ is the positive square-root of the expected mean squared error (expected root-mean-square error) of the estimator.
%\item We are calculating the information gain, uncertainty reduction. Or something proportional to this.
%\end{itemize}

%Although it appears like a restrictive list, what types of experiments does this cover?
%Ramsey interferometry, Mach-Zehnder interferometry. Most experimental demonstrations - B field estimation, frequency estimation, phase estimation, operating regime of LIGO (assuming short measurement time so that the gravitational wave signal does not change much).

\emph{Counting approach:} For a measurement device comprised of $N$ spins, assuming we can prove
\begin{quotation}
Statement 1: A single squeezed state does not provide fundamentally more information on $\theta$, per unit time, than a single spin.
\end{quotation}

It then follows that the information bound on $\theta$ for any $N$ spin ensemble containing $M$ squeezed states is lower than an ensemble of $N$ unentangled spins. The reason being that a squeezed ensemble can be separated into a maximum of $M$ independent state vectors where $M < N$ (this follows from the very definition of entanglement). Assuming these states are independent\footnote{This assumption just ensures that any correlations come solely from entanglement.}, the information provided by the squeezed ensemble is $M$ times that of a single squeezed state and is therefore less than the information provided by the unentangled ensemble. In fact, this counting argument shows that increasing squeezing leads to a worse uncertainty bound since there are less states to average over. Meaning that, if the uncertainty provided by an ensemble containing squeezed states ever surpasses the single particle bound and if the degree of squeezing/entanglement is continuously increased, then at some point the uncertainty must get worse. And the converse: if the measurement uncertainty always improves for arbitrary amounts of squeezing, then the uncertainty never surpasses the single particle bound.

Many mathematical statements equivalent to Statement 1 (excluding the counting argument) have been published (see Appendix\,\ref{App:Proofs}). For example, Wootters showed in 1981 that $\func{I}{\theta,t}$ can be interpreted as distance metric over quantum states \cite{Wootters1981}, meaning that if the path mapped out by unitary evolution is the same for two states then so is the information. This is equivalent to saying that the uncertainty depends only on the sensor response to $\theta$ \cite{Ou1997, Childs2000, Giovannetti2004, Zwierz2012, Pang2014, Yuan2015, Pang2017, Gorecki2020}, therefore states with no enhanced response provide no fundamental sensing advantage. Braunstein, Caves and Milburn provided most of the mathematical content of the proof by showing that for pure states, $\func{I}{\theta,t}$ is given solely by the gradient of the state vector with respect to $\theta$, and does not depend on any intrinsic uncertainty of this state \cite{Braunstein1996}. Here we detail the arguments in full.

\section*{Proof - Time-independent multiplicative Hamiltonian: $\hat{H}(\theta) = \theta \cdot \hat{H}$}
Denote a spin-squeezed state as $\ket{\psi_{SS}(\theta,t)}$, and the state of a single spin as $\ket{\psi_1(\theta,t)}$. Denote the maximum information on $\theta$ that can be provided by any measurement on $\ket{\psi_{SS}(\theta,t)}$, $\ket{\psi_1(\theta,t)}$ as $\funcwrt{I}{\theta, t}{SS}$ and $\funcwrt{I}{\theta, t}{1}$ respectively, then we have:\\

\textbf{Claim:} If $\left(\frac{\mathrm{d}\bra{\psi_{SS}(\theta,t)}}{\mathrm{d} \theta}\right)\left(\frac{\mathrm{d}\ket{\psi_{SS}(\theta,t)}}{\mathrm{d} \theta}\right) = \left(\frac{\mathrm{d}\bra{\psi_1(\theta,t)}}{\mathrm{d} \theta}\right)\left(\frac{\mathrm{d}\ket{\psi_1(\theta,t)}}{\mathrm{d} \theta}\right)$ then $\funcwrt{I}{\theta,t}{SS} \leq \funcwrt{I}{\theta,t}{1}$.\\

\textbf{Physical interpretation:} Squeezed states claim to surpass the SQL by reducing the uncertainty associated with the state, and not by increasing the response of the state to the signal (i.e. the derivative). To refute this claim we need to show that if the gradient of the squeezed state with respect to the signal is the same as that of a single spin, then the information bound on the squeezed state is less than or equal to that of a single spin. As the derivative of a state vector is also a state vector, we can't say one state vector is greater than another, i.e. $\frac{d\ket{\psi_1}}{d\theta} > \frac{d\ket{\psi_2}}{d\theta}$ is not a mathematical statement (the vectors actually exist in different Hilbert spaces). To obtain a non-negative real number, we take the inner-product. Since state vectors are normalised, this operation returns the magnitude of the derivative.\\

\textbf{Proof:} Braunstein, Caves and Milburn showed that for a pure state \cite{Braunstein1996}:
\be
\func{I}{\theta,t} = 4\left[\left(\frac{\mathrm{d}\bra{\psi(\theta,t)}}{\mathrm{d} \theta}\right)\left(\frac{\mathrm{d}\ket{\psi(\theta,t)}}{\mathrm{d} \theta}\right) - \left|\bra{\psi(\theta,t)}\left(\frac{\mathrm{d}\ket{\psi(\theta,t)}}{\mathrm{d} \theta}\right)\right|^2\right].
\label{eq:BCM}
\ee
Working in the Schr\"odinger picture where time-evolution is carried by quantum states, an initial state $\ket{\psi_0}$ evolves in response to the Hamiltonian $\hham(\theta,t) = \theta \cdot \hham$, according to: $\ket{\psi_0} \rightarrow \U(\theta, t)\ket{\psi_0} \equiv \ket{\psi(\theta,t)}$, where $\U(\theta,t) = \func{Exp}{-\iu \theta t \hham /\hbar}$. Writing $\ket{\psi_0}$ in terms of the $K$ eigenstates of $\hham$, denoted by their eigenvalues $\ket{\psi_{E_k}}$ with complex amplitude $\alpha_k$: $\ket{\psi_0} = \sum_{k = 1}^m \alpha_i \ket{\psi_{E_k}}$, we have:
\[ \ket{\psi(\theta,t)} = \sum_{k = 1}^K \func{Exp}{-\iu \theta t E_k/\hbar}\alpha_k \ket{\psi_{E_k}}, \]
where the derivative of this state with respect to $\theta$ is:
\[ \frac{\mathrm{d}}{\mathrm{d}\theta}\ket{\psi(\theta,t)} = \sum_{k = 1}^K (-\iu t E_k/\hbar) \func{Exp}{-\iu \theta t E_k/\hbar}\alpha_i \ket{\psi_{E_k}}.\]
We first derive the maximum information a single spin can provide. For a spin-1/2 with two eigenstates and denoting the eigenvalues of $\hham$ as $\pm \frac{\gamma}{2}$, we have:
\[ \left(\frac{\mathrm{d}\bra{\psi_1(\theta,t)}}{\mathrm{d} \theta}\right)\left(\frac{\mathrm{d}\ket{\psi_1(\theta,t)}}{\mathrm{d} \theta}\right) = \left(\frac{t \gamma}{2\hbar}\right)^2 \left[ |\alpha_1|^2 + |\alpha_2|^2 \right], \]
and
\[ \left|\bra{\psi_1(\theta,t)}\left(\frac{\mathrm{d}\ket{\psi_1(\theta,t)}}{\mathrm{d} \theta}\right)\right|^2 = \left(\frac{t \gamma}{2 \hbar}\right)^2 \left[ |\alpha_1|^2 - |\alpha_2|^2 \right]^2.\]
We can see that $\func{I}{\theta,t}$ is maximised by initially placing the spin in an equal superposition of eigenstates, so that $|\alpha_1|^2 - |\alpha_2|^2 = 0$. Then: 
%\be
\[\funcwrt{I}{\theta,t}{1} = \left(\frac{t \gamma}{\hbar}\right)^2.\]
%\label{eq:single}
%\ee
We have shown that: $\funcwrt{I}{\theta,t}{1} = 4\left(\frac{\mathrm{d}\bra{\psi_1(\theta,t)}}{\mathrm{d} \theta}\right)\left(\frac{\mathrm{d}\ket{\psi_1(\theta,t)}}{\mathrm{d} \theta}\right)$ and reproduced the well-known Heisenberg limit for a single spin \cite{Bollinger1996, Giovannetti2004, Pang2014}. Inserting into Eq.\,\eqref{eq:BCM}, the equality stated in the claim, the maximum information of the squeezed state is:
\[%\be
\funcwrt{I}{\theta,t}{SS} = \funcwrt{I}{\theta,t}{1} - 4\left|\bra{\psi_{SS}(\theta,t)}\left(\frac{\mathrm{d}\ket{\psi_{SS}(\theta,t)}}{\mathrm{d} \theta}\right)\right|^2.
\]%\label{eq:squeeze} \ee
As the second term is a non-negative number, $\funcwrt{I}{\theta,t}{SS} \leq \funcwrt{I}{\theta,t}{1}$ and the proof is complete.

In the above analysis, it may seem like we only consider quantum states and not measurements, however Eq.\,\eqref{eq:BCM} implicitly contains the Born rule. In particular, $\func{I}{\theta,t}$ is a projection of $\frac{\mathrm{d}\ket{\psi(\theta,t)}}{\mathrm{d} \theta}$ meaning that (for this Hamiltonian) the optimal measurement basis is a projection orthogonal to the eigenstates\footnote{Here orthogonal means at an angle of $90^{\circ}$, not anti-parallel.} \cite{Braunstein1996, Childs2000, Giovannetti2004, Pang2017}. Explicitly, for $\hham = \gamma \hat{\sigma}_z/2$ where $\hat{\sigma}_z$ is the Pauli-$z$ matrix, the measurement associated with $\funcwrt{I}{\theta,t}{1}$ is a projective measurement in the $x-y$ plane. Considering just a 2-dimensional space with discrete measurement outcomes, we can denote the measurement results as `1' and `0', thus allowing Eq.\,\eqref{eq:QFI} to be expressed as\footnote{We have used $\func{Pr}{0|\theta,t} = 1 - \func{Pr}{1|\theta,t}$, for a Bernoulli random variable taking only two values.}:
\be
\func{I}{\theta,t} = \left|\frac{\partial \func{Pr}{1|\theta,t}}{\partial \theta}\right|^{2} \frac{1}{\func{Pr}{1|\theta,t}\left(1-\func{Pr}{1|\theta,t}\right)}.
\label{eq:prob}
\ee
Here we can identify $\left|\frac{\partial \func{Pr}{1|\theta,t}}{\partial \theta}\right|$ and $\sqrt{\func{Pr}{1|\theta,t}\left(1-\func{Pr}{1|\theta,t}\right)}$ with the measurement signal and noise in Eq.\,\eqref{eq:SNR}, where the latter is called quantum projection noise \cite{Itano1993}. Note that the description of some quantum states as having intrinsically lower uncertainty is completely absent in this analysis, and in Eq.\,\eqref{eq:prob} the noise and signal are not independent\footnote{We should be careful when equating terms in Eq.\,\eqref{eq:SNR} and Eq.\,\eqref{eq:prob} because they are not identical. In Eq.\,\eqref{eq:SNR} the response of the meter is independent of the measurement, whereas in Eq.\,\eqref{eq:prob}, $\func{Pr}{1|\theta,t}$ depends on the measurement.}.

With the addition of the counting approach detailed earlier, we have proved that ensembles containing pure squeezed states provide a worse uncertainty bound than unentangled ensembles for estimating parameters in time-independent, multiplicative Hamiltonians. In Appendix\,\ref{App:Mixed},\,\ref{App:General} we generalise the proof to the following situations. We show that if we consider:
\begin{itemize}
\item probability mixtures of squeezed states, then the information from these mixed states cannot exceed a single spin.
\item the expected mean information, averaged over the prior probability distribution of $\theta$, then squeezed states cannot outperform a single spin.
\item a modified claim concerning the projection of the gradient $\frac{\mathrm{d} \ket{\psi}}{\mathrm{d}\theta}$ orthogonal to $\ket{\psi}$, then a squeezed state cannot cannot outperform a single spin for estimation of signals in arbitrary time-dependent Hamiltonians $\hham(\theta, t)$.
\end{itemize}

\section*{Discussion}
We have proved a powerful and surprising theorem that allows us to immediately exclude any of the methods proposed in \cite{Caves1981, Walls1983, Wodkiewicz1985, Slusher1985, Wu1986, Slusher1987, Xiao1987, Polzik1992, Wineland1992, Kitagawa1993, Wineland1994, Sanders1995, Kuzmich1998, Soerensen1998, Brif1999, Kuzmich2000, Meyer2001, Orzel2001, Andre2002, Esteve2008, Appel2009, Eberle2010, Gross2010, Koschorreck2010, Koschorreck2010a, Leroux2010, Leroux2010a, LouchetChauvet2010, Riedel2010, SchleierSmith2010, Wasilewski2010, Abadie2011, Luecke2011, Ma2011, Sewell2012, Aasi2013, Taylor2013, Bohnet2014, Muessel2014, Strobel2014, Kruse2016, Polzik2016, Cox2016, Davis2016, Bohnet2016, Hosten2016, Hosten2016a, Linnemann2016, Macri2016, Tse2019, Braverman2019, Schulte2020, Malia2020, Bao2020, PedrozoPenafiel2020, Casacio2021, Gilmore2021, Colombo2022, Greve2022, Malia2022, Eckner2023, Franke2023, Braginsky1974, Braginsky1977, Unruh1979, Braginsky1980, Braginsky1996} from achieving a measurement precision beyond the SQL. Of these works, the following are experimental papers \cite{Xiao1987, Slusher1985, Slusher1987, Polzik1992, Soerensen1998, Kuzmich2000, Meyer2001, Orzel2001, Esteve2008, Appel2009, Eberle2010, Gross2010, Koschorreck2010, Koschorreck2010a, Leroux2010, Leroux2010a, LouchetChauvet2010, Riedel2010, SchleierSmith2010, Wasilewski2010, Abadie2011, Luecke2011, Sewell2012, Aasi2013, Taylor2013, Bohnet2014, Muessel2014, Strobel2014, Bohnet2016, Cox2016, Hosten2016, Hosten2016a, Kruse2016, Linnemann2016, Braverman2019, Tse2019, Bao2020, Malia2020, PedrozoPenafiel2020, Casacio2021, Gilmore2021, Colombo2022, Greve2022, Malia2022, Eckner2023, Franke2023}, the majority of which claim to demonstrate a measurement precision that cannot be obtained without squeezing. %\footnote{We need a method to physically characterise or measure the state dependence on $\theta$ since we proved that if the gradient of a spin squeezed state is the same as that of a single state: $\left|\frac{d\ket{\psi(\theta,t)}_{SS}}{d \theta}\right| = \left|\frac{d\ket{\psi(\theta,t)}_1}{d \theta}\right|$, then both states contain the same amount of information on $\theta$. However the gradient of a state-vector is also a vector, not a non-negative real number, so to get a real number we need to take the inner product.  One might argue, isn't it tautological that if the precision is better, then by definition the gradient was increased! Perhaps people mistakenly attribute this to reduced noise and it is just an explanatory error. The problem is that it is generally obvious and easy to check the state dependence on $\theta$. Just look at the measurement results or the expected state evolution, i.e. the quantum phase accumulate per unit time. It is not the case that the phase change per unit time increases and it is mistakenly reported. Furthermore, if this information is not provided, then \emph{de facto}, sufficient information on the measurement uncertainty has not been presented. As an example take the gravitational wave detector LIGO, where squeezed light is inserted into the interferometer. The imparted phase per photon does not increase when squeezed light is inserted into the interferometer. In Ramsey interferometry with squeezed spins, it is also clear that the phase response is not increased, therefore these techniques cannot beat the SQL. A final point which can be subtle to grasp. I am not suggesting that squeezing did not improve the sensitivity of the LIGO interferometer. The results are quite clear in that regard. I am asserting that squeezing did not improve the sensitivity of the LIGO interferometer beyond a strict impossibility limit for the same number of independent input photons and measurement time.}.
This begs the question:
\begin{quotation}
Why are there so many experimental papers that seem to contradict this proof?
\end{quotation}

The answer is not that our proof is erroneous, but rather that in general, analyses of experiments either do not compare the achieved precision to the SQL or involve fundamental mistakes. Below we outline some common examples in the literature.\\% We acknowledge that the benefits of entanglement in quantum metrology are so ingrained in the consciousness that it can be difficult to overturn this conception. To help we make some ridiculous and outlandish comparisons, not just because it is fun, but because it can help to illustrate how absurd the claims are.\\

%If someone presents you with a machine that produces more output energy than input, a simple appeal to conservation of energy can be used to correct them, likewise proof of the claim, but it is somewhat non-constructive.  but sometimes it is more informative to go into detail and show exactly where the error in logic occurred.

\noindent\emph{Reasons why quantum squeezing has never surpassed the SQL:}%The reason seems to be because most experimental papers use a different definition of the SQL or incorrectly report the measurement precision.

\noindent\textit{A. Comparison to a different limit}
\begin{enumerate}
\item \textbf{Comparison to an imperfect limit:} An improvement in measurement precision for a device operating far away from the SQL is demonstrated \cite{Abadie2011, Aasi2013, Tse2019, Casacio2021}. This is not the same as demonstrating a measurement precision beyond the SQL. % If I promise to improve the performance of your car using \emph{Shacky-mo!}, and provide a demonstration where the speed of a run-down old car is improved slightly, you might wonder whether \emph{Shacky-mo!} will work on your car specifically or improve performance in general. In the context of squeezing, we can ask whether the sensitivity improvement is with respect to a fundamental limit that could not have been achieved otherwise. The answer is no\footnote{Take the work of Casacio et.\,al. \cite{Casacio2021}, where squeezed light lowers the microscope noise without increasing the signal and the authors claim to beat the SQL (they use the term shot-noise limit). In response to Referee 1 (in the peer review file) who notes that the actual performance is underwhelming and far from the state-of-the-art, the authors substantially improve the microscope performance by aligning the microscope. If the performance can be improved through alignment it is clear that comparison is not to a strict impossibility limit. In the context of LIGO \cite{Aasi2013, Tse2019} the analogy of a run-down car is less apt, since this is one of the most sensitive detectors that has ever been developed. However, one can clearly read in LIGO's 2022 \href{https://dcc-llo.ligo.org/public/0182/G2200736/001/UpdateonLVKDetectors.pdf}{update} that significant technical noise sources remain which mean that the interferometer is operating far away from the SQL. A factor of 2 -- 3 sensitivity improvement for a device operating orders of magnitude away from the SQL is not the same as demonstrating an improvement beyond this limit.}.

\item \textbf{Tautological definition of the measurement precision:} A different quantum limit is introduced. This limit is defined as the precision of a measurement device \emph{as it is and without modifications}. Modifications that utilise entanglement are allowed, whilst forbidding any other changes to the device \cite{Wu1986, Xiao1987, Taylor2013}.% on Of course if we forbid anyone changing anything in the measurement device then the precision cannot be improved. The tricky part is to allow modifications that utilise entanglement but forbid any other changes to the device. One example of this definition is the work of \cite{Wu1986, Xiao1987}. 
%\item \textbf{Time not fully accounted for:}
\end{enumerate}
Note, in both of these cases, if the device was already operating at the SQL, and squeezing implemented whilst keeping the number of particles and time constant, then the sensitivity cannot improve.

\noindent\textit{B. Analytic errors}
\begin{enumerate}
\item \textbf{Analysis of measurement noise not uncertainty:} Conflation of noise and uncertainty is one of the most common analytical mistakes in spin squeezing, see e.\,g. \cite{Orzel2001, Esteve2008, Appel2009, Gross2010, Koschorreck2010, Koschorreck2010a, Leroux2010, LouchetChauvet2010, Riedel2010, SchleierSmith2010, Sewell2012, Cox2016, Hosten2016, Kruse2016, Colombo2022, Greve2022, Malia2022, Franke2023}. As an example, assume a projective measurement in the $x$-basis is performed on a single spin pointing in an unknown direction in the $x-y$ plane. Denote this state $\ket{a}$ and label the outcomes of the observable $\hat{S}_x = \hat{\sigma}_x/2$ as $\{0, 1\}$. The uncertainty in estimating the $x$-component of the spin direction is actually independent of the measurement variance. If ten experiments produce the following data-set: $\{1,1,1,1,1,1,1,1,1,1\}$, the SQL has not been surpassed, even though the variance is less than 1/4. Compare to the data-set obtained from a different state $\ket{b}$: \{0,1,1,1,0,0,1,1,1,0\}. Someone with knowledge of $\ket{a}$, $\ket{b}$ can prove this is not a statistically anomaly by showing $\left(\Delta \hat{S}_x\right)^2_a < \left(\Delta \hat{S}_x\right)^2_b$. Here $\left(\Delta \hat{S}_x\right)^2$ is the variance of the $\hat{S}_x$ operator and  the subscript indicates the state that the variance is calculated with respect to. Despite this mathematical inequality, our uncertainty in estimating the $x$-component of $\ket{b}$ is the same as for $\ket{a}$. Ultimately we do not want measurements with less \emph{noise}, but measurements that provide lower \emph{uncertainty} on the spin direction\footnote{In the literature $(\Delta \myvec{S})^2$ is often incorrectly equated with the variance of some spin operator $(\Delta \hat{A})^2$. These are not the same quantities! In general a collective spin observable $\hat{\bm{S}}_z = \sum_{i=1}^N \hat{S}_z$ is measured, such that a variance $\left(\Delta \hat{\bm{S}}_z\right)^2 <|\myvec{S}|/4$ is presented evidence of an enhanced precision, where $|\myvec{S}|$ denotes the spin vector length. Note, in Eq.\,\eqref{eq:SNR}, $\Delta \myvec{S}$ is the uncertainty in estimating the spin direction, it is not the (square-root) variance of an operator, nor is it the square-root variance of some dataset.}. Often prior information from state preparation (and the value of $\theta$) is used to ensure that the measurement statistics have low variance without reducing the uncertainty in estimating $\theta$.
%It is easy to perform measurements with low noise -- throw away the data. Given $N$ identical copies of a state, the noise can be reduced by a factor of $\sqrt{N}$ by measuring only one state. According to squeezing logic, this strategy surpasses the SQL. Thus we can see that some processes which increase the measurement noise actually reduce the measurement uncertainty.

\item \textbf{Replace the signal gradient with the contrast:} The measurement signal (c.f. the derivative in Eq.\,\eqref{eq:prob}), is replaced with a constant term that does not depend on $\theta$ or the measurement basis. A common error when the measurement precision is defined, is to replace the gradient with the spin vector length, characterised by the `contrast', `visibility', `coherence' or `Wineland squeezing parameter' \cite{Meyer2001, Esteve2008, Gross2010, Leroux2010, Leroux2010a, Riedel2010, Hosten2016, PedrozoPenafiel2020, Malia2020, Malia2022, Eckner2023}. Again a measurement basis is chosen to reduce the measurement noise without accounting for the commensurate reduction in signal. I.e. the experiment is designed so that $\func{Pr}{X|\theta,t} \rightarrow 0$, but the loss of signal $\frac{\partial \func{Pr}{X|\theta,t}}{\partial \theta} \rightarrow 0$ is not fully accounted for, and is often assumed to be constant. To some extent, the operation regime of Advanced LIGO \cite{Aasi2015, Tse2022} and QND proposals \cite{Braginsky1980, Kimble2001} suffer from this error.

%\item \textbf{Incorrect definition of the SQL:} The SQL is commonly defined a time-independent measurement precision or defined with a square-root time dependence, e.g. \cite{Leroux2010a}.
\end{enumerate}

%What is going on here? Obviously I need to address that in at least some experiments, squeezing does improve the measurement precision, albeit not beyond the SQL. At the very least one should admit there is a dilemma here which should be investigated. When discussing the assumptions, an explanation for some of these observations is provided in Appendix\,\ref{App:Assumptions}).

There are valid questions on how one should interpret this proof. For instance it only addresses fundamental quantum measurement noise\footnote{It also neglects measurement back-action onto the signal. However back-action is minimised when using a single particle sensor.}, other technical sources of experimental noise should be reduced. One question is to what extent this is possible, or more broadly whether the information bound of Eq.\,\eqref{eq:BCM} is achievable? This may allow for some practical regime where squeezing can be of advantage. However, as we note in Appendix\,\ref{App:Assumptions} invoking this explanation is tantamount to saying quantum mechanics is not correct. Perhaps a more nuanced interpretation is to note that no experiment can reach the precision bounds set by quantum mechanics since reaching these bounds requires perfect, instantaneous measurements with no decoherence. But you can't have it both ways. Claims that squeezing can surpass the SQL do not present that message, indeed quite the opposite.

%\textit{There are valid objections to this proof. Only deals with fundamental quantum measurement noise, for other sources of course the noise should be reduced. But entanglement doesn't help with this. Perhaps there is a reason why we can't physically reach this bound for a single spin, and a much stricter limit holds for measurements of a single spin. but we can for an entangled spin. Note for a single spin the optimal measurement is also relatively easy to implement, it is just a projective measurement in the basis spanned by the superposition states. of the spin into one of the eigenstates. Perhaps, decoherence, ...}

\section*{Conclusion}
At a fundamental level, quantum mechanics dictates that measurement uncertainty is a statistical sampling phenomenon. Some statistical analyses of measurement precision mistakenly associate the fundamental randomness of measurement observations as a noise which should be reduced. Here we have shown that approaches to improve the measurement uncertainty by reducing this noise cannot succeed. To misquote Rolf Landauer \cite{Landauer1998} -- \textit{The signal is the noise!}

By relating metrological performance to the number of separable states we have proved that squeezed ensembles cannot outperform unentangled ensembles in sensing. The proof was inspired by the proposition of a general quantum mechanical uncertainty limit -- the amount of information per unit time that can be extracted from a non-separable state is bounded by $I\left[E,t\right] < (t/\hbar)^2$, where $E$ denotes the energy of any state with respect to another \cite{McGuinness2021a}. Here wavefunctions are considered the fundamental information unit in computers and sensors, not particles. My position is that this information per unit time limit holds for all entangling procedures, not just squeezing, and its application extends beyond sensing to all tasks in quantum information including computation. Finally, it should be apparent that in the field of quantum physics, the peer-review process has failed (see also \cite{McGuinness2022, McGuinness2023, McGuinness2023a}) and this has led to a loss of scientific progress.

\bibliographystyle{naturemag}
%\bibliographystyle{unsrt}
\bibliography{}%{references.bib}%{}

\newpage

\section*{Appendix}
\subsection{Photon interferometry and vacuum noise}\label{App:LIGO}
%In the following analysis I am going to be making several approximations, or assumptions, simplifications. The purpose of these simplifications is to allow me to address the simple question ``Can squeezing provide a fundamental improvement over non-squeezing". To do this, we want to remove all other noise sources. For example, advanced ligo used an additional LO field to produce power variations that are detected by the photodiode.... Equivalent to lock-in detection for improving the signal to noise of many devices.
%
In his seminal paper on quantum mechanical noise in an interferometer, Caves noted a fundamental noise source that ``can be attributed to vacuum (zero-point) fluctuations in the electromagnetic field, which enter the interferometer from the unused input port" \cite{Caves1981}. One typical physical description of this noise is that when an electromagnetic wave enters the interferometer at the first beam-splitter, vacuum fluctuations entering the unused port cause an unequal splitting of the electromagnetic field. Caves noted an equivalent view that treats photons as particles and attributes this imbalance to the ``random scattering of the input photons at the beam-splitter". As a result more of the input light is directed to one arm of the interferometer and less into the other, in turn affecting the output light intensity at both ports. Yet another description is that photon shot-noise randomises the arrival time of photons on the detector \cite{Tse2019}. It is commonly understood that this noise cannot be removed, except through inserting squeezed vacuum into the unused input port to reduce these fluctuations. Here I present a different physical description of the interferometer, one that changes the emphasis of the noise and I believe, helps better predict experimental results.

As in the main text, one can learn a lot by reducing the experiment down to its individual components and providing a physical description of this basic experiment. So what happens when a single photon enters the interferometer? Using a particle description of the photon, which randomly scatters at the beam-splitter, we encounter a problem. The photon only traverses one arm, thus there is no interference. No matter the path difference, the photon exits the unused port with 50\% probability. This is not what is observed experimentally. Furthermore, the idea of each photon traversing just one arm of the interferometer leads one to infer that the output signal depends on interference between \textit{separate} photons. This does not capture the basics, rather it is interference of a photon with \textit{itself} that is relevant\footnote{More precisely it is interference between basis states of a single wavefunction, which can include more than one photon for entangled states.}. I also find the suggestion that vacuum fluctuations somehow prevent one from creating an equal superposition state hard to accept -- the creation of equal superposition states is a basic concept in quantum metrology and quantum information processing. Importantly, this description also incorrectly predicts experimental results. For example, the observed probability (up to experimental imperfections) to detect a photon at the unused port is given by $p = \sin^2(\phi/2)$, where $\phi$ is the relative phase difference. If vacuum fluctuations always enter the unused port to randomise the observed photon statistics, then how can one explain this observation? Take $\phi = 0$. Without inserting any squeezed vacuum into the unused port, one can predict the experimental outcome with remarkable accuracy. Fringe visibilities in excess of 90\% have been demonstrated so far for single photon interference. This should not be possible if vacuum fluctuations of order $\sqrt{N}$ are a real and fundamental source of noise.

A different description to those presented so far is to assume the outgoing photon wavefunction at the beam-splitter is a perfect superposition of basis states (as far as technical tolerances allow), so that a relative phase difference between returning basis states determines their complex coefficients\footnote{Here basis states are eigenstates of photon number in each mode of the interferometer.}. A photodetector, placed at the unused port, performs a measurement of the photon wavefunction, and detects a photon with probability given by the squared amplitude of the state in this mode. In this description, there is no noise in the interferometer which we need to remove. The noise is a statistical effect, due to probabilistic measurement and is the result of the Born rule. When multiple independent photons are inserted into the interferometer, the basis states of each individual wavefunction interfere upon return, and the summation of many individual probabilistic measurements yield the detected signal. We do not need to worry about balancing the light intensity in each arm of the interferometer, except to manufacture an ideal beam-splitter. Note, in this description, if the interferometer is operated so that the unused port is dark, then although laser power noise and quantum projection noise is reduced \cite{Itano1993}, so is the signal since $\frac{\partial p}{\partial \phi} \rightarrow 0$.% Operating close to a dark fringe provides no fundamental benefit if there is no technical noise, however it can have practical advantages. If readout noise at the photodetector dominates, then operating half-way between a dark and bright fringe improves sensitivity, but if laser power noise dominates, then it is advantageous to operate close to a dark fringe. \textit{But not directly on the dark fringe as the sensitivity goes to zero. Need to check this}\footnote{What about collective or non-linear effects occurring when multiple identical photons enter at the same time? Well, these generally occur when identical photons enter at different input ports, e.g. Hong-Ou-Mandel. We don't need to worry about these because we assume nothing enters the input port.}.

To convincingly demonstrate that vacuum fluctuations are not a fundamental noise source in photon interferometry, we can go one step further and remove the beam-splitter, thereby preventing vacuum fluctuations from entering the interferometer. %To do this, we need to temporally separate photons in one arm of the interferometer from photons in the other arm, and we need a device that measures the phase of wavefunctions. 
To illustrate the idea, take photons produced from a pulsed laser, and replace the beam-splitter with a flippable (double-sided) mirror, starting flipped down so that incident photons enter the longitudinal arm of the interferometer, after which the mirror is flipped up, so that the second pulse enters the transverse arm. The mirror remains in this position so that photons returning from the longitudinal arm are sent out the antisymmetric (unused) port, after which it is flipped down allowing photons returning from the transverse arm to output at the same port. Now we do want to compare the relative phase between separated optical pulses. % If we can replace the photodiode normally positioned at the antisymmetric output of the interferometer with a phase sensitive detector, this would allow us to compare the phases of both optical pulses and detect a shift in one arm of the detector.
However, as optical photodiodes are sensitive only to the light intensity and not the phase, this mode of operation seems forbidden. It turns out that coherent ensembles of spins can be used as phase sensitive detectors.

As a first approximation, let's treat the photodetector quantum mechanically and the optical pulses as classical control fields. In this description, gravitational waves change the relative phase of control fields acting on spin qubits. The first (longitudinal) optical pulse performs a $\pi/2$-rotation on the ensemble, where the optical phase is stored coherently in the spin direction of the ensemble. This phase can be `readout' with the second (transverse) optical pulse which performs a $\pi/2$-rotation around an axis determined by the relative phase of the two pulses (see \cite{McGuinness2021a} for a demonstration with microwave photons). Finally the spin state along-$z$ should be projectively readout. For a single spin, the probability to remain in the initial $\ket{0}$ state is: $p(\ket{0}) \approx \sin^2(\phi/2)$, where $\phi$ now refers to the phase difference between the two pulses. If the spin energy transition is perfectly resonant to the laser frequency, time-delays between the pulses have no effect on the final spin state. %In fact we do not even need an interferometer \emph{per se}, to detect gravitational waves. We can remove the mirror and make use of just the longitudinal arm of the interferometer. Now the phase of returning photons is simply compared to the phase of photons outputted from the laser directly which are either simultaneously or sequentially applied to an ensemble of spins. This mode of operation only allows detection of gravitational waves along one axis and requires that the laser has good phase stability.

An advantage of treating spins rather than photons, as the measuring device, is that it is much easier to perform quantum non-demolition (QND) measurements on spins. Not only that, sensitivity to power fluctuations can be removed by operating near $\phi = \pi$ so that $p(\ket{0}) \approx 1$ and the spins are in a `dark fringe'. Consider fluctuations slow enough so that both the longitudinal and transverse pulse have the same intensity. Now instead of $\pi/2$-rotations, the spins are rotated through an angle $\Theta$. The $\pi$-phase shift of the second pulse ensures the spins are brought back to $\ket{0}$. There is a fundamental physical principle behind this observation, and it allows quadratic suppression of all amplitude fluctuations. The photon wavefunction is mapped to the spin state, so that conjugate observables of the photon wavefunction (phase and amplitude) are mapped to conjugate spin observables ($\hat{S}_z$ and $\hat{S}_x$). Measurements that obtain information on the photon phase cannot also reveal information on, and are therefore insensitive to the photon amplitude.

To see this, assume we only want to detect gravitational waves along the longitudinal axis, which corresponds to measuring the phase $\varphi$ of the first optical pulse. The first pulse is mapped to the spin state resulting in:
\[\ket{\psi} = \cos(\Theta)\ket{0}+e^{-\iu \varphi}\sin(\Theta)\ket{1}.\]
The pulse amplitude (proportional to $\Theta$) is mapped to $S_z$ and (the cosine of) $\varphi$ mapped to $S_x$. As phase and amplitude are conjugate observables, a $\hat{S}_z$ measurement gives no information on $\varphi$ -- one observes $p(\ket{0}) = \cos^2(\Theta)$, $p(\ket{1}) = \sin^2(\Theta)$. Measurement of the laser phase $\varphi$, requires measurement of a conjugate observable. In order to perform such a measurement, we do not need to send an optical pulse into the transverse arm of the interferometer, we just need to implement a $\hat{S}_x$ (or $\hat{S}_y$) measurement. Although there are technical challenges in performing such a perfect measurement, the postulates of quantum mechanics, assume that that it is possible. Furthermore, we do not need to analyse the back-action this measurement induces on the mirror and there is no radiation pressure trade-off. We have now described a setup that allows realisation of the QND type measurements originally proposed by Braginsky, Vorontsov and Thorne \cite{Braginsky1980} and more recently by Kimble et.\,al. \cite{Kimble2001}. In the language of uncertainty squeezing and QND measurements, one would say that the meter remains in an eigenstate, and measurements of this state have minimum uncertainty. However, when the measurement basis is chosen to make the readout noise zero, the meter response is quadratic to phase shifts and is insensitive to small phase shifts.

Finally, one should again consider the case, where just a single photon enters the interferometer and is detected with just a single spin. The spin-photon interaction should be tailored so that the combined system is in an entangled superposition: $\frac{1}{\sqrt{2}}\left(\ket{0_S}\ket{1_P}+e^{-\iu \varphi}\ket{1_S}\ket{0_P}\right)$, where the $S$, $P$ subscripts denote the spin and photon states respectively. In this scenario there is the outstanding issue on defining the phase of a single photon. I argue that $\varphi$ is a physical property of the experiment, and that it can be observed by interfering this state with a suitable (single) photon state.

Although further details are required to make this discussion rigorous, that does not make it useless. All physical descriptions involve some approximations, the key is to find which ones are critical to obtaining accurate predictions. One benefit of this perspective is to demonstrate that vacuum fluctuations can be removed from an interferometer. We have reduced the problem of gravitational wave detection to estimating spin direction, a problem which is explicitly addressed in the main text. Of course, vacuum fluctuations are replaced with something else -- quantum projection noise when the spin state is readout -- but as we have shown, this noise cannot be removed through squeezing. Furthermore, this shifts the focus, we should focus on reducing technical noise -- laser power fluctuations, beam-splitter, detector efficiency. Most importantly, there is an achievable parameter space where the predictions of this analysis diverge greatly from treatments of vacuum fluctuations. For low input power into LIGO where number of photons is small so radiation pressure is very small, we can define the minimum possible sensitivity per unit time that one can achieve with $N$ \textit{input} photons into the interferometer. If we now implement squeezing and use the same number of \text{input} photons into the interferometer and the same measurement time, we can ask whether the interferometer sensitivity can surpass this limit? Our analysis show that the answer to this question is no.\\

%One might ask ``What is going on in LIGO experiments then?" ``How can they report sensitivity improvements with squeezing when the above analysis says it is impossible?" I will start by noting that it is remarkably difficult to find a clear and understandable description of the physical process in LIGO that accurately describes how squeezing improves the sensitivity of the spectrometer. Secondly, as far as I am aware, the sensitivity improvements of Quantum-Enhanced Advanced LIGO -- where squeezed light is inserted into the detector -- with respect to Advanced LIGO are not reported in terms of more detection events with higher signal-to-noise ratio, instead they are presented in terms of noise spectra (see this 2022 LIGO \href{https://dcc-llo.ligo.org/public/0182/G2200736/001/UpdateonLVKDetectors.pdf}{update}). I contend that the issue with this analysis is that LIGO operates near a dark-fringe meaning that the output photons are highly post-selected. In fact, photons that experience a significant phase shift are much more likely to be detected. Benchmarking the sensitivity in terms of the output number of photons makes it seem like LIGO has high sensitivity, but if we look at the total number of photons inserted into the interferometer, LIGO is operating many orders of magnitude away from the SQL.

%How accurate is this analysis, where you assume no photons entering the unused port? Isn't this unphysical? 

%Under this perspective, when $N$ photons of perfectly squeezed light are incident on the beam-splitter, they all distribute to the same arm of the interferometer and exit at the same port. Importantly the phase shift experienced by these photons is the same as for a single photon. In effect the squeezed photons collectively act as a single photon, so that the photon statistics and information on the phase shift is identical to what would be observed by passing a single photon through the interferometer.

%It is not lost on me that this proof seems to contradict fundamental descriptions of interferometer performance. This section tries to reduce the intellectual work-load in translating the results proved in the context of spins to the context of light interferometry. Also some of the cognitive dissonance involved when comparing the proof to such descriptions.\\

%In order to make any mathematical description of a physical experiment understandable and concise certain assumptions need to be made. The same holds for noise analysis. Here, for the single particle uncertainty limit to be attained we need to make certain assumptions. These assumptions also apply to the entangled case, and they do not affect the inequality that we have proved, but it is worthwhile to discuss them. Making the assumptions also has a benefit of allowing us to precisely locate where certain noise terms arise and then discuss whether approaches exist to remove the noise. We are talking about a precision limit per unit time and per number of photons. LIGO SQL is sometimes defined as precision limit for any possible number of photons. Depends on the test mass. There are many other assumptions that LIGO makes, for example low losses or no losses. Thermal black-body radiation from the equipment etc. They are all fine because they are small with respect to other noises, and maybe also there is some sense in which they are not fundamental, there are ways that we can reduce them, by cooling, better design etc. But the other ones cannot be reduced. It's also worthwhile to make a special point of interest on where people disagree. I agree with the community, that when the laser power increases the radiation noise will increase, and this will limit the sensitivity. I disagree with the LIGO community that squeezed light can reduce the photon shot noise below some unentangled fundamental limit. That is the part I would like to focus on.

%\textit{Need to point out that we are really only considering photon shot noise, not radiation-pressure noise. In analysis of LIGO these are inversely proportional to each other. Increase one, reduce the other. But also depends on mass of the test mass mirrors. Even if we consider radiation-pressure noise, we can show that squeezing is worse. The reason is that squeezed states have more radiation pressure than a single spin. Braginsky et. al. 2002/2003 paper show that we can consider the quantum back-action noise by treating the test masses quantum mechanically, where measurement of the position influences their momentum and vice versa. The measurement is performed by the photons. Or we can disregard this noise and instead focus on the laser light. The issue with this idea is that we cannot just disregard this noise and set it to zero. If we want to set it to zero, then we need to increase the test masses infinitely. If we do this, then the radiation-pressure noise has no impact. Also, the SQL in most descriptions for spins, the signal is treated classically. There is no back action of the spins and their measurement on the value of $\theta$. It's not so much do we treat the test masses as quantum mechanical or classical, it's more do we allow interaction of the photons with the test masses to change their position and momentum.}

%\textit{I think that in this version of squeezing, we are reducing the photon shot noise at the expense of increasing the radiation pressure noise. Now one arm of the interferometer, experiences all of the radiation pressure, whereas the other arm experiences none of the pressure. We could make a different version of squeezing where when $N$ photons of perfectly squeezed light are incident on the beam-splitter, they evenly distribute to the each arm of the interferometer. The question is if they exit at the same port or not. If they do, then they do not produce any radiation pressure noise. One part of Cave's argument that we can disregard fluctuations in the mean power of the laser - because mean radiation-pressure forces on both masses do not affect the measurement (of the difference), is that it does not consider the final output signal, and this depends on the mean power of the laser and thus this noise does impact the sensitivity. But maybe they are using balanced quadrature detection which removes this noise.}

\subsection{A confusing history of early ideas to surpass the SQL}\label{App:History}
The idea of using special projective measurements (not states), with reduced noise originated in discussions on improving the sensitivity of gravitational wave detectors \cite{Braginsky1974, Braginsky1977, Unruh1979, Braginsky1996}. These low noise measurements, called quantum non-demolition (QND) measurements, are performed by choosing a measurement operator so that the meter is in an eigenstate of the observable and the variance of the measurement observable goes to zero \cite{Unruh1978, Braginsky1980}. However no entanglement is required to observe such a variance since separable eigenstates also yield zero variance results. One unfortunate relic of this early analysis which persists today is that the SQL is often defined as the minimum obtainable variance on one observable when the variance equals that of another conjugate observable. Often this statement is incorrectly reduced to saying that the uncertainty in measuring spin-direction is equivalent to the spin noise, giving the SQL: $\Delta \myvec{S} = \sqrt{|S|}/2$. Thus when the spin noise is $< \sqrt{|S|}/2$, the SQL for measuring spin direction has been surpassed. Initially the idea was to take advantage of better measurements, not entangled states, however around the time of Caves' proposal for squeezed vacuum fluctuations in the context of photon interferometry \cite{Caves1981}, emphasis on states of the meter with reduced uncertainty started. These methods were then extended to spin states in Ramsey interferometers, called spin-squeezed states \cite{Kitagawa1991, Wineland1992, Kitagawa1993}.

We can summarize some reasons why QND measurements do not surpass the SQL. First, the requirement of perfect prior knowledge about the state of the meter in order to perform a QND measurement is self-defeating\footnote{If we are being pedantic, one can learn about \emph{which} eigenstate the meter is in, and perform discrimination.}. Put another way, if we want the noise to remain low, the meter must remain in an eigenstate, but if the meter simply remains in an eigenstate then nothing is happening?! Also there is an inherent lack of recognition of the difference between `uncertainty' and `noise'. In general, QND schemes neglected to consider that the measurement signal and measurement noise are not independent, measurements with reduced noise -- where the meter is near an eigenstate of the measurement operator -- also have less signal. In fact, Wootters showed in 1981 \cite{Wootters1981}, that even in a $2^N$ dimensional Hilbert space the optimal measurement yields a probability function of the form: $\func{Pr}{1|\theta} = \func{cos^2}{\frac{m}{2}(\theta - \theta_0)}$, where $m$ is an integer and $\theta_0$ is a constant; thus when the probability to measure an observable goes to one or zero, the gradient disappears (Unruh used related arguments in response initial proposals \cite{Unruh1978}). Moreso, Itano et. al. noted that since projective measurements are not perfect, it is often preferable to choose a measurement that maximises projection noise \cite{Itano1993}. The message here is that simply reducing measurement noise does not guarantee better precision, and in general measurements which maximise noise are better. Finally, Giovannetti, Lloyd and Maccone proved, for a restricted set of Hamiltonians, that entangled states and not measurements are required to surpass the SQL \cite{Giovannetti2006}.

Another relic of early suggestions to use QND measurements to surpass the SQL is their current\footnote{It seems that my characterisation of QND measurements as a discounted historical idea is not accurate. Carlton Caves informs me that LIGO are still working on QND measurements as outlined here \cite{Kimble2001}.} use in the preparation of squeezed spin states \cite{Kuzmich2000, Appel2009,  Koschorreck2010, LouchetChauvet2010, SchleierSmith2010, Sewell2012, Bohnet2014, Cox2016, Hosten2016, Hosten2016a, Braverman2019, Malia2020, Greve2022, Malia2022}. The idea is to place an ensemble of spins initialised in $\ket{0}^{\otimes N}$ in an optical cavity, where a $\pi/2$-pulse rotates $\myvec{S}$ into the $x - y$ plane. A QND measurement is then performed on the ensemble to project a sub-ensemble of the spins to either $\ket{0}$ or $\ket{1}$. This measurement is generally a dispersive population measurement of the spins, where an optical probe pulse is inserted into the cavity and the phase (frequency) shift of the output light is recorded. The shift is proportional to the number of atoms projected into $\ket{0},\ket{1}$. Apparently, after performing this QND measurement, the ensemble is in an entangled state with less noise and therefore improved measurement precision.

I strongly disagree with this description. After the QND measurement, I would describe a single state in the ensemble as being in a probability mixture of $\ket{0}, \ket{1}$ and $\ket{x}$\footnote{Actually, not exactly $\ket{x}$, the electromagnetic field of the optical probe slightly rotates this state around $z$.}. To refute this description at least one piece of information is required: a definition of the phase relationship between entangled basis states. For example, one can write an entangled Bell state as: $(\ket{00} + \ket{11})/\sqrt{2}$, and although it is generally overlooked, a well-defined phase between the states $\ket{00}$ and $\ket{11}$ is required to generate and make use of this entangled state (in this example the state can be written $(\ket{00} + e^{\iu \phi }\ket{11})/\sqrt{2}$ with $\phi = 0$). Otherwise we have a mixed state and the best we can say is that for any measurement we will record either a `0' or a `1', each with 50\% probability. None of the papers that use QND preparation of a ``squeezed" spin state define or make use of any phase relationship between the entangled basis states. It is true that knowledge of the phase of the initial $\pi/2$-pulse and the intensity the probe light is required when measuring these states, but these parameters do not determine $\phi$ (see Assumptions and FAQ's sections below).

\subsection{Proof assumptions}\label{App:Assumptions}
For completeness, we list some assumptions of the proof below. In particular we assume:
\begin{itemize}
\item The postulates of quantum mechanics.
\item The entire dependence on $\theta$ is encoded in $\hham(\theta,t)$. The starting state $\ket{\psi_0}$ and the measurement operator are independent of $\theta$.
\item Apart from the initial state vectors that evolve to $\ket{\psi_{SS}(\theta,t)}, \ket{\psi_1(\theta,t)}$, everything else in the external universe is the same when comparing unentangled and entangled ensembles. The prior information on $\theta$ and the ability to apply control fields is the same. In fact, my understanding is that the squeezing community claim even unitary evolution is the same for these two states, therefore the only difference of consequence is the starting state (and the measurement).
\item Each copy of the same state is independent and measurements on copies of the same state are identical. This assumption ensures that any enhancement over the SQL comes solely from entanglement, not other correlations, additional information or different control on copies.
\end{itemize}
\textbf{Comment on Assumption 1:}
I discuss one assumption of the analysis presented here and use a simple example to illustrate how it is critical to the proof. The assumption is the (Born rule) postulate of quantum mechanics -- the probability of obtaining a measurement outcome $m$ when measuring a state $\ket{\psi}$ is equal to $p(m) = \bra{\psi}\hat{M}^{\dagger}_m \hat{M}_m\ket{\psi}$, where $\{\hat{M}_m\}$ is the collection of self-adjoint measurement operators associated with outcomes $m$ (Postulate 3 in \cite{Nielsen2000}). This postulate does not account for the fact that measurements are never perfect in experiments and that some quantum states are easier to measure.

Take fluorescent readout of the spin states where the state $\ket{0}$ fluoresces strongly and $\ket{1}$ does not fluoresce. If we compare readout of the following spin states, performed by scattering photons and counting the total number of photons detected (a collective spin measurement):
\begin{align*}
\mathrm{State}\,1: \quad \frac{1}{\sqrt{2}}(\ket{0}& + \ket{1})\\
\mathrm{State}\,2: \;\; \frac{1}{\sqrt{2}}(\ket{00}& + \ket{11})
\end{align*}
then the second state is easier to measure. For readout along $z$, the photon difference for measurements of the second state is twice that when reading out the first state\footnote{One could argue that these measurement results have greater variance, i.e. the variance of the operator $(\Delta S_z)^2$ is greater for the second state. Again, this operator variance argument is often incorrectly employed to say that measurement observables of one state provide more information on the spin direction than the other. Apart from the fact that we are comparing two different operators here, at a fundamental level in terms of the amount of information provided, I disagree. If the measurements are perfect, then both observables provide the same information on the spin direction. Also the measurement variance is arbitrary to some extent. We can take readout of the first state and multiply by an arbitrary large number to increase the variance (or multiply by a small number to reduce the variance). In terms of measurement statistics, this does not change the information contained in the measurements.}. Whenever state readout is not perfect, then (under the fair sampling assumption) the second state provides more information. A few points of emphasis. While this is a practical limitation in every real experiment, the proposition that it represents a fundamental precision limit in quantum squeezing is equivalent to saying quantum mechanics is incorrect. Also, increasing the readout fidelity removes this loophole since the information difference reduces as the readout improves. Importantly, in spin squeezing literature, the second state (or its extension to $N$ particle Hilbert space) is not being used, the following is used instead:
\[ \frac{1}{\sqrt{2(1+\epsilon^2)}}\left(\epsilon\ket{00} + \left(\ket{10}+\ket{01}\right) + \epsilon\ket{11}\right) \]
where $\epsilon$ is real, $\epsilon <1$ for a squeezed state, $\epsilon \ll 1$ in the limit of strong squeezing, and $\ket{10}$, $\ket{01}$ cannot be distinguished by a collective spin measurement along $z$ \cite{Kitagawa1993, Ma2011, Andre2002, Vuletic2023} (see FAQ's section).\\


\noindent\textbf{Assumptions for physical interpretations:}

The following assumptions are not required to derive the mathematical inequalities, but if one then uses those inequalities to claim that the precision of a squeezed device cannot surpass the precision of an ideal unentangled ensemble, that argument requires further assumptions. I.e. when going from the maths to physical interpretations of the proof one makes implicit use of these assumptions. 
\begin{itemize}
\item In general (quantum) measuring devices are not initialised in $\ket{\psi_0}$, therefore we should account for the time and resources required to prepare this state. Using $\ket{\psi_{0,1}}$, $\ket{\psi_{0,SS}}$ to differentiate between $\ket{\psi_0}$ for a single spin and a squeezed state respectively, then we assume the time and resources to generate $\ket{\psi_{0,1}}$ is not greater than $\ket{\psi_{0,SS}}$. In every proposal and experiment that I am aware of, at least one copy of $\ket{\psi_{0,1}}$ is used to create $\ket{\psi_{0,SS}}$. For example, a $\pi/2$-rotation is performed on at least one spin as the first step in creating entangled spin states. For Ramsey interferometry, the spin is in $\ket{\psi_{0,1}}$ after the $\pi/2$-rotation, however additional gates are required to produce $\ket{\psi_{0,SS}}$. To create squeezed photon states, single photons are inserted into a non-linear optical element, or the interferometer together with an additional field. Thus the assumption that, in general it is more technically demanding and time-consuming to generate the ideal $\ket{\psi_{0,SS}}$ compared to the ideal $\ket{\psi_{0,1}}$, seems correct. In fact, when generating entangled NOON, GHZ and CAT type states, (I claim) it precisely because the converse of this assumption does not hold that these entangled states do not outperform a single particle\footnote{I.e. for NOON, GHZ and CAT states to reach the Heisenberg limit or outperform a single spin, the time and resources needed to generate $\ket{\psi_0}$ for these entangled states is assumed to be short.}. I would like to remove this assumption, but have been unable to find a rigorous proof. One approach is to compare the distance between the initialised basis state $\ket{0}^{\otimes N}$ and $\ket{\psi_{0,1}}$ or $\ket{\psi_{0,SS}}$. However, $\func{Dist}{\ket{0},\ket{\psi_{0,1}}} \nless \func{Dist}{\ket{0}^{\otimes N},\ket{\psi_{0,SS}}}$ for any distance metric I have found.
\item A single spin produces no greater measurement back-action onto the signal compared to a single squeezed state. Again it seems like this assumption can be made rigorous by proving that a single spin produces minimal back-action for the amount of information extracted, and showing this is not the case for a single squeezed state. We can remove this assumption with respect to LIGO, since a single photon has lower radiation pressure than a squeezed state. Our treatment of the SQL assumes that measurement back-action is zero. I.e. we treat the signal classically and assume it has a fixed value that does not vary. This is standard in theoretical analyses of the Heisenberg limit, especially in metrology with spins. For LIGO, this is the low light power operation regime. Again we emphasise that as squeezed states produce more back-action than a single photon our results apply to the high power regime.
\item In going from (Fisher) information to measurement uncertainty, we make assumptions on the estimator obtained from measurements of single spin compared to measurements of a squeezed state. Note, we do not need to assume that the minimal variance unbiased estimator exists \cite{Kay1993}, which may not always be the case. For readout of $\ket{\psi_1(\theta,t)}$ we just need the existence of any estimator with lower mean squared error than the best estimator obtained from readout of $\ket{\psi_{SS}(\theta,t)}$. In practice, this is easy to check, since state readout of a single spin only gives two outcomes and finding the estimator is straightforward.
\item We assume that the information limit for ideal measurements of a single spin can be reached. More precisely, if the bound is not strict, meaning there is a gap between theory and practice, that squeezed states cannot fit in this gap. In general, if we define $t$ as the total measurement time, the finite $\pi/2$-rotation time required to create $\ket{\psi_0}$, and the finite readout time prevent $\funcwrt{I}{\theta,t}{1}$ from being reached. However this is also the case for squeezed states. We also assume no other technical noise or decoherence, or more precisely that measurements of single particle states are less susceptible to this noise. Importantly, this assumption allows us to define $\funcwrt{I}{\theta,t}{1}$, and we note these assumptions are generally used in derivations of the SQL and $\Delta \est{\theta}_1$. Going beyond this assumption seems to require a modification of the postulates of quantum mechanics. It is a position I strongly advocate.
\\
\end{itemize}

\subsection{Alternative proofs of a noise independent uncertainty limit}\label{App:Proofs}
%\textit{From my perspective, the number of available works presenting equivalent content to the proof presented here, highlights how big the blind-spot in the physics community is. Indeed (as the mathematics is well-known) it seems like the main difficulty is simply entertaining the idea that entanglement could lead to a worse measurement uncertainty, or having the conviction to follow the mathematics despite being told otherwise.} 
The following gives some sketches for alternative approaches to obtaining the same proof as in the main text. They are not rigorous and more work is required to close technical details. 
\begin{enumerate}
\item We could make the observation that most experimental implementations of squeezing such as quadrature detection in optical interferometers or spin readout in Ramsey interferometers result in a binary outcome measurement. Thus we can use the expression given in Eq.\,\eqref{eq:prob}:
\[\func{I}{\theta,t} = \left|\frac{\partial \func{Pr}{1|\theta,t}}{\partial \theta}\right|^{2}\frac{1}{\func{Pr}{1|\theta,t}\left(1-\func{Pr}{1|\theta,t}\right)}.\]
It is left as an exercise to show that any modifications to $\func{Pr}{1|\theta,t}$ that increase $\func{I}{\theta,t}$ necessarily increase $\left|\frac{\partial \func{Pr}{1|\theta,t}}{\partial \theta}\right|$. More involved is to show that this is the case even for probability distributions taking $2^N$ discrete values, thus we do not need to restrict ourselves to a two-dimensional probability space.

\item We can follow Wootters \cite{Wootters1981} and show that $\func{I}{\theta,t}$ is a distance metric on state vectors, characterising the (angular) distance between the start and final state. As this distance is parameterized by the state evolution it follows that if state evolution is the same for two states, then the information is also the same. Equivalently, the information provided by a path of states is given by the gradient vector. We need to use the fact that state vectors are normalised, so they have the same length.

\item In Refs. \cite{Wootters1981, Braunstein1994, Childs2000, Giovannetti2006, Jones2010, Pang2017} amongst others, an uncertainty bound depending only on the maximal difference in eigenvalues of $\frac{\partial \hham(\theta,t)}{\partial \theta}$ is provided. Minimising the uncertainty requires placing the meter in an equal superposition of states with extremal eigenvalues, $\mu_{\mathrm{max}}(t), \; \mu_{\mathrm{min}}(t)$:
\[
\ket{\psi(t)} = \frac{1}{\sqrt{2}}\left(\ket{\mu_{\mathrm{max}}(t)} + e^{\iu \Phi}\ket{\mu_{\mathrm{min}}(t)}\right),
\]
where $\Phi$ is an arbitrary phase. The optimal state is isomorphic to a single spin, since it is just a 2-dimensional space and the measurement uncertainty per unit time is limited by the maximum response speed of the meter to changes in $\theta$.

\item We could observe that the uncertainty on $\theta$ in Eq.\,\eqref{eq:SNR} is neatly divided into two time-zones (this is in contrast to Eq.\,\eqref{eq:prob}). The first contribution is during sensor interaction with the signal and the second occurs after interaction and is just readout of the sensor. Now the benefit from squeezing closely resembles early proposals to surpass the SQL using QND measurements with lower uncertainty (see Appendix\,\ref{App:History}). But Giovannetti, Lloyd and Maccone proved that entanglement cannot provide any benefit at the measurement stage for a restricted class of Hamiltonians \cite{Giovannetti2006}. For a rigorous proof we also need to show that the two terms in Eq.\,\eqref{eq:SNR} are independent. For example, one might say that putting the whole ensemble in a NOON state reduces the ability to measure the direction of the meter -- as compared to $N$ copies of the same state, thus states that interact faster also happen to be states that cannot be measured well. However the uncertainty in measuring the direction of a single NOON state is the same as measuring the direction of a single particle state, the worse uncertainty simply comes from having less states to average over. The temporal separation achieves this to the extent required, since after interaction, the first term is fixed and we do not yet need to make any decision on the measurement.

\item We could apply results from quantum state tomography, where the uncertainty in estimating the spin direction of an unknown state just depends on the number of copies of the state \cite{ODonnell2016, Aaronson2018}.

\item We could observe the hint provided by looking at state discrimination, i.e. when we know that $\theta$ can take only one of two discrete values. Childs, Renes and Preskill \cite{Childs2000} showed that driving the state to one of two orthogonal states minimises the discrimination uncertainty and the best one can do is to drive as quickly as possible. There is no benefit to reducing readout noise here since it already goes to zero.\\
\end{enumerate}

\subsection{Proof - Average expected information and mixed states}\label{App:Mixed}
The key to these extended proofs is that we are dealing with probability distributions over non-negative real numbers, so we don't have to worry about any complex number magic. In comparing summations of non-negative real numbers where each term in one sum is less than each corresponding term in another sum, by convexity we have an inequality.\\

\noindent \textbf{Mixed states:} To prove the claim when the squeezed state is a mixed state $\ket{\rho_{0,SS}}$, we can write the initial state as probability distribution over $L$ pure states, with probabilities $p_l$: $\ket{\rho_{0,SS}} = \sum_{l=1}^L p_l \ket{\psi_{0,SS}}_l$. As we have assumed that the initial state does not depend on $\theta$, we can use the property that the Fisher information is convex over mixed states when $p_l$ is independent of $\theta$ \cite{Liu2020}.\\

%Then as the information from any individual pure state does not exceed a single spin, the average information from the mixed state cannot either. This is equivalent to  Although, here we do not need to assume independence of probabilities because the optimal state independent of $p_l$.

%For mixed states, when the probabilities are parameter independent, i.e. when $\func{Pr}{\theta} = \mathrm{Pr}$, then the quantum Fisher information is convex:
%\[ \func{I}{\func{Pr}{\cdot}\rho_1(\theta)+(1-\func{Pr}{\cdot})\rho_2(\theta),H} \leq \func{Pr}{\cdot}\func{I}{\rho_1(\theta),H}+(1-\func{Pr}{\cdot})\func{I}{\rho_2(\theta),H}\]


\noindent\textbf{Average expected information:} Eq.\,\eqref{eq:QFI} just provides the information around a given value of $\theta$. The expected information, averaged over the prior probability distribution of $\theta$ is:
\be \langle \func{I}{\theta,t}\rangle = \int^{\Theta} \mathrm{d}\Theta p(\theta=\Theta) \func{I}{\Theta,t}, \quad \langle \func{I}{\theta,t}\rangle = \sum_{j=1}^J p(\theta = \Theta_j) \func{I}{\Theta_j,t}
\ee
where $p(\theta = \Theta)$ is the prior probability that $\theta$ takes on some value $\Theta$. The LHS expression is the expected information when the prior probability is continuous and the RHS when $\theta$ can take $J$ discrete values.

Given that the information on any value of $\theta$ is not greater than a single spin, we can show that this holds for the expected information averaged over the probability distribution of $\theta$. First we note that the optimal state is independent of the value of $\theta$ and therefore also independent of the prior probability of $\theta$. It is not however the case that the optimal measurement is independent of prior probability distribution of $\theta$. To additionally prove that no measurement can produce more information, we use the result of Giovannetti, Lloyd and Maccone who showed that entanglement does not provide any advantage at the measurement stage \cite{Giovannetti2006}. Thus we only need to consider the amount of information encoded in the state (before measurement).\\

\subsection{Proof - Arbitrary time-dependent Hamiltonian: $H(\theta, t)$}\label{App:General}
We motivate the proof by noting that if we could solve the unitary evolution of a single spin and show that it remains in the optimal state for all evolution times, then this would imply $\left|\bra{\psi_1(\theta,t)}\left(\frac{\mathrm{d}\ket{\psi_1(\theta,t)}}{\mathrm{d} \theta}\right)\right|^2 = 0$. However, without a general method of solving time-evolution, even when restricting ourselves to a single spin, we cannot make that statement. An equivalent approach is to use the information bound proved by Pang and Jordan for general time-dependent Hamiltonians \cite{Pang2017}:
\be
\func{I}{\theta,t} \leq \left| \int^t_{0} \mathrm{d}t' \left(\mu_{\mathrm{max}}(t') - \mu_{\mathrm{min}}(t')\right)/\hbar \,\right|^{2},
\label{eq:Pang}
\ee
where $\mu_{\mathrm{max}}(t), \; \mu_{\mathrm{min}}(t)$ are the maximum, minimum eigenvalues of $\frac{\partial \hham(\theta,t)}{\partial \theta}$. Eq.\,\eqref{eq:Pang} expresses a simple concept, the information per unit time is limited by the maximum response speed of the meter to changes in $\theta$. This quantity is characterised by the eigenvalues of $\frac{\partial \hham(\theta,t)}{\partial \theta}$ (c.f. how the eigenvalues of $\hham$ determine the evolution speed of any quantum state). Minimising the uncertainty requires placing the meter in an equal superposition of states with extremal eigenvalues, $\mu_{\mathrm{max}}(t), \; \mu_{\mathrm{min}}(t)$:
\[
\ket{\psi^{\mathrm{Opt}}(t)} = \frac{1}{\sqrt{2}}\left(\ket{\mu_{\mathrm{max}}(t)} + e^{\iu \Phi}\ket{\mu_{\mathrm{min}}(t)}\right),
\]
where $\Phi$ is an arbitrary phase \cite{Braunstein1996, Giovannetti2006, Childs2000}.

If we could solve the unitary evolution under $\hham(\theta, t)$, show that $\ket{\psi_1(\theta,t)} = \ket{\psi^{\mathrm{Opt}}(t)}$, and find a projective measurement that saturates the bound, then we can saturate Eq.\,\eqref{eq:Pang} and show that $\funcwrt{I}{\theta,t}{1} = 4\left(\frac{\mathrm{d}\bra{\psi_1(\theta,t)}}{\mathrm{d} \theta}\right)\left(\frac{\mathrm{d}\ket{\psi_1(\theta,t)}}{\mathrm{d} \theta}\right)$. Unfortunately in general this is not possible: for example when both the eigenvalues and eigenstates of $\hham$ depend on $\theta$. Therefore we need to modify our claim to address not the gradient of $\ket{\psi}$, but rather the projection of $\frac{\mathrm{d}\ket{\psi}}{\mathrm{d} \theta}$ orthogonal to $\ket{\psi}$: $\left(\frac{\mathrm{d}\ket{\psi}}{\mathrm{d} \theta}\right)_{\perp} \equiv \frac{\mathrm{d}\ket{\psi}}{\mathrm{d}\theta}-\ket{\psi}\bra{\psi}\left(\frac{\mathrm{d}\ket{\psi}}{\mathrm{d}\theta}\right)$.\\

\textbf{Claim:} If $\left(\frac{\mathrm{d}\bra{\psi_{SS}(\theta,t)}}{\mathrm{d} \theta}\right)_{\perp}\left(\frac{\mathrm{d}\ket{\psi_{SS}(\theta,t)}}{\mathrm{d} \theta}\right)_{\perp} = \left(\frac{\mathrm{d}\bra{\psi_1(\theta,t)}}{\mathrm{d} \theta}\right)_{\perp}\left(\frac{\mathrm{d}\ket{\psi_1(\theta,t)}}{\mathrm{d} \theta}\right)_{\perp}$ then $\funcwrt{I}{\theta,t}{SS} = \funcwrt{I}{\theta,t}{1}$.\\

\textbf{Proof:} As noted by Braunstein, Caves and Milburn \cite{Braunstein1996}
\[ \left(\frac{\mathrm{d}\bra{\psi(\theta,t)}}{\mathrm{d} \theta}\right)_{\perp}\left(\frac{\mathrm{d}\ket{\psi(\theta,t)}}{\mathrm{d} \theta}\right)_{\perp} = \left(\frac{\mathrm{d}\bra{\psi(\theta,t)}}{\mathrm{d} \theta}\right)\left(\frac{\mathrm{d}\ket{\psi(\theta,t)}}{\mathrm{d} \theta}\right) - \left|\bra{\psi(\theta,t)}\left(\frac{\mathrm{d}\ket{\psi(\theta,t)}}{\mathrm{d} \theta}\right)\right|^2 \]
and Eq.\,\eqref{eq:BCM} proves the claim. This statement is equivalent to looking at the data-set $\{X_i|\theta,t\}$ resulting from projective measurements of $\ket{\psi_{SS}(\theta,t)}$. If this data-set shows no greater dependence on $\theta$ than a single spin, then it cannot contain fundamentally more information on $\theta$ than projective measurements of a single spin.\\

%\subsection*{Philosophical question}
%Who should claim the proof? Clearly it was shown by Braunstein, Caves and Milburn. This is more a discussion on what constitutes a proof. In mathematics, given some axioms, it is generally fixed what is provable or not provable in that axiomatic system. So one can say that the person writing the axioms has proved everything it is just a crank the handle type operation. But in general, some problems are described to require some special insight in order to solve them. The person providing that insight is usually credited with the proof, and sometimes the subsequent people who provide limited insight but do the arduous crank-the-handle work in extending to other situations or rigorously checking and defining the proof are also credited. But how do we quantify insight? It is ill-defined and a grey area. One person, could claim that it is obvious to them and they proved it, but did not provide the work, in other cases this is not so.

%For a mathematical journal, field it would be obvious that BCM did the mathematical proof. However in the past 30 years since this work, it is also obvious that they (and subsequent people) did not recognize the implications of this proof. Grigori Perelman 


\subsection{FAQs}
\noindent \textbf{I find it difficult to accept your statement that the whole field is wrong without more evidence. Could you help point out some specific errors in papers.}

I have started doing that for many experiments, see here \cite{McGuinness2021, McGuinness2022, McGuinness2023, McGuinness2023a}. Let's go through an explicit example for spin-squeezing. All you need to do is ask some simple questions about what is going on and see if the answers make sense. Some questions that might help:
\begin{quotation}
In spin squeezing, what entangled state is being created?
\end{quotation}
It is surprisingly difficult to find a direct answer to this question. In place of a mathematical definition of the squeezed state, a pictorial representation on the Bloch sphere is often used. For a system consisting of two particles, Mark Kasevich and Vladan Vuleti\'c inform me that the state that everyone is preparing is \cite{Vuletic2023}:
\be \frac{1}{\sqrt{2(1+\epsilon^2)}}\left(\epsilon\textcolor{green!50!black}{\ket{00}} + \left(\textcolor{red!50!black}{\ket{10}+\ket{01}}\right) + \epsilon\textcolor{green!50!black}{\ket{11}}\right)
\label{eq:squeezed}
\ee
where $\epsilon$ is real, $\epsilon <1$ for a squeezed state, $\epsilon \ll 1$ in the limit of strong squeezing, and $\ket{10}$, $\ket{01}$ cannot be distinguished by a collective spin measurement along $z$ (I have added colours for emphasis). Equivalent definitions and the extension to $N$ particles are given in [\textcolor{blue!50!black}{14},\,\textcolor{blue!50!black}{43},\,\textcolor{blue!50!black}{97}].

Let's assume this is the input state to a Ramsey interferometer, where a relative phase between basis states accumulates and is compared to the phase of the final readout pulse in the interferometer. The most important point to highlight about such a highly squeezed state ($\epsilon \ll 1$) is that a measurement of collective spin $\hat{\bm{S}}_z = \sum_{i=1}^N \hat{S}_z$ gives no information about relative phase. It is a metrologically useless state because one always obtains the same measurement result $\bm{S}_z = 0$ regardless of any phase shifts. There are several other important characteristics of this state which are in complete discrepancy with what is presented in the literature.
\begin{itemize}
\item The green basis states accumulate a relative phase twice that of a single particle. If the $N$ particle squeezed state has components: $\textcolor{green!50!black}{\ket{0}^{\otimes N}}$ and $\textcolor{green!50!black}{\ket{1}^{\otimes N}}$, they accumulate a relative phase, $N$ times that of a single particle.
\item The relative phase accumulated by the red basis states is not observable by a collective spin measurement along $z$. In fact, a perfect measurement of these basis states always gives the same observable $S_z = 0$, independent of the phase of the final readout pulse.
\item For $\epsilon = 1$, although one might say the state in Eq.\,\eqref{eq:squeezed} is not squeezed, it is still entangled. On average, we would expect to observe Ramsey fringes with a different contrast and period compared to the same number of independent spins. For a two-spin ensemble, if the unentangled ensemble has a period of $2 \pi$ and the fringe contrast goes from $1$ to $-1$, then the entangled ensemble should have a period of $\pi$  and fringes going from $1/2$ to $-1/2$.
\item When squeezing is increased, i.e. $\epsilon \rightarrow 0$, we would expect to see contrast of the Ramsey fringes disappear when performing a collective spin measurement. It is true that the measurement variance is reduced, since we always obtain the same result $S_z = 0$, but this gives no information on the phase of the final readout pulse. Think about what is happening here, squeezing just reduces the amplitude of the Ramsey fringes, without affecting the period. How can this improve the measurement precision?
\end{itemize}

Again, points to emphasise. Measurements of this state give completely different data, compared to measurement of the same number of unentangled spins. In particular, the \textit{period} is reduced and the contrast is reduced by the same amount. Also, when squeezing is increased, although we observe that the measurement variance reduces it is clear that this leads to \textit{worse} phase estimation.

If we accept that the squeezed state defined above is actually being created, the next question one could ask is:
\begin{quotation}
What experimental evidence is presented, to verify the creation of this state?
\end{quotation}

Note, in the following papers, a collective spin measurement along $z$ is performed on the spin squeezed state \cite{Appel2009, Gross2010, Leroux2010, Leroux2010a, LouchetChauvet2010, SchleierSmith2010, Bohnet2014, Muessel2014, Strobel2014, Cox2016, Hosten2016, Linnemann2016, Braverman2019, Malia2020, Bao2020, PedrozoPenafiel2020, Greve2022, Malia2022, Riedel2010, Hosten2016a, Colombo2022}. No evidence is presented of a state with the above characteristics (except that the Ramsey fringe contrast reduces and the measurement variance reduces). Most strikingly, when a plot of the measurement observable vs. readout phase is presented, the period is the same as for the unentangled ensemble, see:
\begin{enumerate}
\item Figure\,2 of \cite{Gross2010}
\item Figure\,2 of \cite{Leroux2010}. Data is only presented for the variance as a function of rotation angle.
\item Figure\,5 of \cite{LouchetChauvet2010}
\item Supplementary Figure\,A6 of \cite{SchleierSmith2010}
\item Figure\,2 and Figure\,3c of \cite{Bohnet2014}
\item Figure\,1c and Figure\,2 of \cite{Muessel2014}. No data is presented for the unsqueezed ensemble, but the expected dependence on the readout phase is the same as presented for the squeezed ensemble (a period of $2\pi$).
\item Figure\,4A of \cite{Strobel2014}. This data corresponds to measurements on a state, which is prepared as a spin squeezed state and evolved for 25\,ms. After 25\,ms squeezing is lost, although entanglement remains?! In Figure\,2, data is presented for squeezing as a function of tomography angle. 
\item Figure\,3(d) of \cite{Cox2016}. A plot of spin noise as a function of the phase $\psi$ of the final readout pulse.
\item Figure\,2 of \cite{Linnemann2016} where a $2\pi$ phase dependence is observed. See also Figure\,3(a) inset.
\item Figure\,2 of \cite{Braverman2019}, data is presented for the variance as a function of rotation angle.
\item Figure\,2c of \cite{PedrozoPenafiel2020}, a plot of variance as a function of tomography angle is presented for the spin squeezed state
\item Figure\,2a of \cite{Riedel2010}, the same angular dependence is observed for squeezed and unsqueezed ensembles.
\item Figure\,4c of \cite{Colombo2022} where the phase response for a squeezed and unsqueezed ensemble are directly compared. See also Figure\,4b for a plot of spectroscopic enhancement as a function of rotation angle for the spin squeezed state.
\item Figure\,2a (inset) of \cite{Eckner2023} where the phase response for a squeezed and unsqueezed ensemble are superposed.
\item Figure\,4a of \cite{Franke2023} where the phase response for a squeezed and unsqueezed ensemble are plotted.
\end{enumerate}
Note for measurements of noise or variance, twice the angular or phase dependence is observed compared to population measurements (i.e. half the period). The same is also observed for measurements of unentangled ensembles when the correct rotation direction is chosen. Sometimes the sensor response is pictorially represented, and no experimental data is presented, see: 
\begin{enumerate}
\item Figure\,1 of \cite{Appel2009}
\item Figure\,2 of \cite{Leroux2010a}
\end{enumerate}
In fact, these two points: no mathematical definition of the squeezed state, a spin response the same as an unentangled state, are what forced me to take the approach I did in the main text.

What should be abundantly clear, is that the squeezed state defined in Eq.\,\eqref{eq:squeezed} (or the $N$ particle equivalent) is not created. So, the natural question to ask is
\begin{quotation}
What state is being created?
\end{quotation}
Summarizing the experimental evidence that has been presented. Collective spin measurements of ``squeezed" ensembles have the same phase dependence as unentangled ensembles and with lower contrast. In addition, the measurement noise (spin variance) is observed to reduce for particular readout angles. An important piece of information that the above papers do not mention, is that the measurement noise (spin variance) also reduces for particular readout angles of unentangled ensembles (see for example Figure\,11 of \cite{Itano1993}). Despite many papers claiming the contrary, it is not correct that reduced noise for particular projective measurements is evidence of squeezing or entanglement. Exactly the same dependence is seen for measurements on unentangled systems.

Finally, we can read the description of how the ``squeezed" state is created. In the following papers, a ``squeezed" state is prepared by performing a QND measurement on the ensemble, i.e. a projective measurement of the spin population along $z$ \cite{Appel2009, Leroux2010, Leroux2010a, LouchetChauvet2010, SchleierSmith2010, Bohnet2014, Cox2016, Hosten2016, Braverman2019, Malia2020, Bao2020, PedrozoPenafiel2020, Greve2022, Malia2022, Hosten2016a, Colombo2022}. In general, this measurement is weak, in that it only projects a small proportion of spins to the $\ket{0}, \ket{1}$ eigenstates of $z$. After performing the QND, the authors state that the ensemble is now in an entangled ``squeezed" state.

Based on the evidence presented in these papers, this is definitely not the description I would use. I would describe some of the spins as being in $\ket{0}$ or $\ket{1}$ conditional on the QND measurement result. That is basic quantum mechanics, after a projective measurement, the system is in an eigenstate of the measurement observable. The spins that are not measured, are unaffected except to experience an a.c. Stark shift from the QND measurement. So, this is my explanation of what happening in the above papers. The spins are prepared in an unentangled state: $\ket{x}^{\otimes N}$. Some of these spins are projected to $\ket{0}$, $\ket{1}$. The resulting ensemble can be described as one where most of the spins are initialised to $\ket{x}$ and some are in $\ket{0}$, $\ket{1}$. This ensemble is then used to perform Ramsey interferometry. The resulting contrast is worse than an ensemble initialised to $\ket{x}^{\otimes N}$ and the phase response is the same. As far as I am aware, this is the only description that fits to all of the evidence presented and is theoretically consistent.\\

\noindent\textbf{Why do you restrict your analysis to (Fisher) information, why not analyse the actual measurement uncertainty? Are you hiding something behind this definition?}

I analyse the information and not the measurement uncertainty for a reason that is both somewhat uninteresting and at the same time points to a bigger issue in quantum metrology. The reason is that it is easy to violate bounds on the measurement uncertainty, whereas the same is not true for the information. This makes the uncertainty bound non-rigorous.

In general, the following inequality is used to relate the information (from a single measurement) to the estimation uncertainty \cite{Braunstein1996, Giovannetti2004, Giovannetti2006, DemkowiczDobrzanski2012, Pang2014, Yuan2015, Pang2017}:
\be
\langle (\Delta \est{\theta})^2\rangle \ge \frac{1}{\func{I}{\theta,t}},
\label{eq:info}
\ee
where $\langle (\Delta \est{\theta})^2\rangle$ is the expected mean squared error of the estimator. Thus the uncertainty $\Delta \est{\theta}$ is the positive square-root of the expected mean squared error of the estimator.

For $t = 0$, no information is obtained on the signal $\theta$, thus one would expect the uncertainty to be infinite. However if we have any amount of prior information on $\theta$, then the uncertainty bound is violated, because at $t = 0$ the uncertainty is not infinite, it is given by our prior uncertainty on $\theta$. To account for prior information we could reformulate our uncertainty definition and work in terms of relative uncertainty reduction. As Eq.\,\eqref{eq:info} relates the amount of information provided by a measurement, the inverse square-root determines the relative uncertainty reduction with respect to the prior uncertainty. However, even using this definition of relative uncertainty reduction it is possible to violate the bound. Take the following discrete prior probability distribution of $\theta$:
\[ p(\theta = \pi/2) = 1/2; \quad \quad p(\theta = \pi) = 1/2, \]
i.e. $\theta$ can take one of two values with equal probability. For this prior, the bound on relative uncertainty reduction given by Eq.\,\eqref{eq:info} can be violated by an arbitrary amount. To see this, note that the initial uncertainty in $\theta$ is $\pi/4$. A measurement is performed where a single spin sensor is driven to one of two orthogonal states (e.g. $\ket{0}$, $\ket{1}$) dependent on the value of $\theta$; so that if $\theta = \pi/2$, the spin ends up in $\ket{0}$ and if $\theta = \pi$, the spin ends up in $\ket{1}$. If the readout of the spin is perfect, then regardless of how long it takes to perform unitary evolution, the final uncertainty on $\theta$ is zero. As $\func{I}{\theta,t}$ is in general at most a quadratically increasing function of $t$ and is always finite, we can obtain a violation. Even allowing for non-prefect readout as long as the probability to obtain a correct result is $>1/2$, the measurement can be repeated thus producing an exponential reduction in uncertainty and violation of the bound for arbitrary numbers of measurements. Similar arguments have been observed in obtaining violations of the Heisenberg time-energy uncertainty relation \cite{Aharonov2002} and used as evidence that quantum computers can exponentially outperform classical computers \cite{Atia2017}.

These examples show that if we really want to obtain a rigorous bound on the measurement uncertainty, we need to be careful in considering the prior information, because the final uncertainty depends on both the information provided from the measurement and our prior uncertainty. Mathematically, the reason why we can violate Eq.\,\eqref{eq:info}, known as the Cram\'er-Rao bound, in the above example is because the posterior probability distribution does not satisfy the regularity condition 
\[ \8{\frac{\partial \mathrm{ln}\left(\func{Pr}{X|\theta,t}\right)}{\partial \theta}} = 0 \quad \mathrm{for \; all}\;\theta\]
where the expectation is taken with respect to $\func{Pr}{X|\theta,t}$ (see section on Cram\'er-Rao lower bound in \cite{Kay1993}). This condition is assumed in the derivation of the uncertainty bound on the mean squared error of the estimator. If we restrict ourselves to prior probability distributions that ensure the posterior probability distribution satisfies the regularity condition, then we can prove that the bound holds.

Performing the above analysis adds complexity and forces us to put restrictions on the prior probability distribution to avoid pathological priors. This also makes the analysis less general. Most importantly this analysis is lengthy, mathematically involved, distracting and does not address the issue at stake. We don't care about the prior probability distribution and how that impacts the final measurement uncertainty, what we want to know is the following
\begin{quotation}
Do measurements of entangled squeezed states give more information on a signal than measurements of unentangled states?
\end{quotation}
Perhaps when we convert from information to uncertainty, the result may violate the bound of Eq.\,\eqref{eq:info}, but we can say that if the answer is no, then measurements of entangled squeezed states will not provide a lower uncertainty than measurements unentangled states (see also the discussion in Appendix.\,\ref{App:Assumptions}).\\

\noindent\textbf{What about the Heisenberg uncertainty relations for conjugate observables. Don't they disprove this work. Are you saying they are wrong?}

Instead of operating in the Schr\"odinger picture we could analyse how unitary evolution acts as a mapping of quantum states in the Heisenberg picture. Then, for $\hham(\theta, t) = \theta \hham$, we have: $\frac{\mathrm{d} \U(\theta,t)\ket{\psi_0}}{\mathrm{d}\theta} = -\iu t \hham \, \U(\theta, t)/\hbar \ket{\psi_0}$. As $\U$ and $\hham$ commute, expressing $\ket{\psi_0}$ in terms of the eigenstates of $\hham$, we have:
%\begin*{align}
\[
\begin{split}
& \left(\frac{\mathrm{d}\bra{\psi(\theta,t)}}{\mathrm{d} \theta}\right)\left(\frac{\mathrm{d}\ket{\psi(\theta,t)}}{\mathrm{d} \theta}\right) = \bra{\psi_0}\left(-\iu t \hham \, \U(\theta,t)/\hbar\right)^{\dagger}\left(-\iu t \hham \, \U(\theta,t)/\hbar\right)\ket{\psi_0} \\
& = (t/\hbar)^2 \sum_{k = 1}^K \bra{\psi_{E_k}} \left(\func{Exp}{\iu \theta t E_k/\hbar}E_k \alpha_i^{\ast}\right)\left(\alpha_i E_k\func{Exp}{-\iu \theta t E_k/\hbar}\right) \ket{\psi_{E_k}}\\
& = (t/\hbar)^2 \sum_{k = 1}^K  E_k^2 |\alpha_i|^2,\\
\end{split}
\]
and
\[ \left|\bra{\psi(\theta,t)}\left(\frac{\mathrm{d}\ket{\psi(\theta,t)}}{\mathrm{d} \theta}\right)\right|^2 = \left|-\iu t/\hbar \bra{\psi_0}\U^{\dagger} \hham \,\U\ket{\psi_0}\right|^2 = (t/\hbar)^2\left(\sum_{k = 1}^K  E_k |\alpha_i|^2\right)^2.\]
%\end*{align}
Thus, for this Hamiltonian, maximising information on $\theta$ means maximising the variance of $\hham$\footnote{In general we want to maximise the variance of the generator of translation with respect to $\theta$ (see \cite{Braunstein1996, Pang2014, Pang2017}). For a time-independent Hamiltonian, if $\U$ and $\partial \hham/\partial \theta$ commute, then we want to maximise the variance of $\partial \hham/\partial \theta$.}: $(\Delta \hham)^2 \equiv \langle \hham^2 \rangle - \langle \hham \rangle^2$ with respect to $\ket{\psi_0}$. For unitary evolution this also means maximising the variance of $\hham$ on the output state. We have:
\[ \func{I}{\theta} = \frac{4t^2}{\hbar^2}\left(\Delta \hham\right)^2.\]

Due to the correspondence between different pictures, increasing $\left(\Delta \hham\right)^2$ in this context only comes about through increasing the state response to $\theta$. In this analysis, it is incorrect to equate $\left(\Delta \hham\right)^2$ as the uncertainty in estimating the energy eigenvalues or the Hamiltonian $\hham$. It is the variance of the measurement results, assuming perfect measurements. Importantly, we see that states with large variance in measurement outcomes (i.e. noise) give the most information on $\theta$. This is in direct contradiction with the analysis provided by the squeezing community.

There are some technical aspects to this analysis which means that working in the Schr\"odinger picture is preferable. In the main text, we assume that two state vectors have the same dependence on $\theta$. However in the Heisenberg picture, the dependence is included in the operator not the state vector, so we cannot make that assumption. Also when comparing a single spin with a squeezed state, we are comparing the variance of two different operators with different dimension. We don't have that issue in the Schr\"odinger picture because we are comparing inner products on normalised vectors. Even when the vector spaces are different dimensions, this is still ok. On a related point, defining the measurement variance is a tricky endeavour because processing of raw experimental data is always required to obtain meaningful estimates. Experimentally one does not measure energy by mathematically applying the $\hham$ operator to a system upon which the outcome $E_k$ is returned. Instead, most often in spin experiments, spins are readout fluorescently where photons are collected and differences in photon counts are related to measurement outcomes. However, there is nothing to stop us from multiplying the photon data by a large number and artificially increasing the variance, or for that matter either not collecting data or performing disingenuous subtractions to reduce the variance. These operations do not improve our measurement precision and it can be hard to identify such artificial enhancements and separate them from actual improvements. 

Take the following phase estimation experiments. In one experiment, the state to be measured is $\left(\ket{1} + e^{\iu \varphi}\ket{0}\right)/\sqrt{2}$. Compare to another experiment with entangled Bell state $\left(\ket{1 1}+ e^{\iu \varphi}\ket{0 0}\right)/\sqrt{2}$. What are the uncertainty bounds on estimating $\varphi$ from these two states? The measurement outcomes of the second experiment produce twice as many photons as the first. But in the limit of perfect fidelity readout, these extra photons do not provide more information on $\varphi$ -- they are correlated! Squeezing people might claim that because the measurement observable from the Bell state has twice the variance of the single spin state, then it has a worse uncertainty. Furthermore by measuring a conjugate observable we would obtain a better uncertainty.

In fact, assuming the postulates of quantum mechanics are correct, both states provide fundamentally the same amount of information on $\varphi$. The variance argument presented in spin squeezing is incorrect in many respects. For one, it is a simple rescaling error. The analytic mistake introduced by looking at variances of observables becomes even harder to identify when we start talking about measurements on ensembles where the contribution of each state gets blurred out. It is much easier to pick up on this error when comparing individual normalised states in the Schr\"odinger picture (see the comment on Assumption 1 in Appendix.\,\ref{App:Assumptions}). 
\\


%\textbf{Are squeezed states just NOON, GHZ or CAT type states? Maybe people are using a different description of the same physics.}
%
%No. The squeezing community
%
%It is difficult to find a description of the entangled squeezed state in the literature. In these papers you will find a pictorial representation of the quantum state in place of a mathematical description, where a Gaussian blur on a (Bloch) sphere is replaced with an elongated blur on a (Bloch) sphere \cite{Gross2012, Pezze2018}. I personally find this representation very hard to follow, because any state on the Bloch sphere is isomorphic to a single spin, so I would represent it by an arrow in the spin direction, not a blur. Indeed this lack of a definition of the squeezed state is what led me to take the approach in the main text. I had to rely on how the authors describe their experiment and present the supporting data, rather than explicitly analysing the sensitivity of the state vector that is used.
%
%I am not aware of any experimental paper in spin squeezing where


%\textbf{Why restrict the analysis to just a single non-separable state vector, why not analyse the whole ensemble?}

%Because this method is cleaner. Differentiating between the signal, noise is often confusing and blurry. For example, when measuring two copies of a state, does the signal stay the same and the noise go down, or does the noise go up and the signal also go up? By how much. The main reason is that single state vectors are normalised. 

%In this case, because everything is static, then we just need to use a normalisation argument. There is nothing particularly special about the number of photons here, we have to assign a given level to a data-value. We could also consider two different single spins, one that fluoresces brightly and one less bright. Yes, there are technical and practical reasons why it is better to chose the bright spin, but it is not correct to say that the bright spin achieves an information limit that is impossible with the darker one. As long as we are close to high fidelity, then the difference in values is even less important.

%\textbf{What is the connection between photons and spins?}
%
%There are mathematical analyses showing photon interferometry is equivalent the interferometry with spins. Thus the same applies for photons. But just to say they are mathemetically equivalent is not to say that they are physically identical or that the experiments are the same. Two distinct differences are - it is hard to do non-demolition experiments with photons. Measuring a photon state generally destroys the photon. Secondly, beam-splitters are replaced with $\pi/2$-pulses, so vacuum fluctuations do not enter into spin experiments. What are they replaced with? Quantum projection noise.\\


%In your paper, you use the argument that the HL is not strict as a reason why entanglement doesn't improve in general. Aren't you just making the same assumption here. Isn't that hypocritical, you are just committing exactly the same error you accuse others of.




%\textbf{Why do you say that peer-review has failed. Why don't you do something about it?}
%I have been trying to. I have contacted many of the authors of experimental papers. I have written a review of experimental papers. I have directly written to the journal editors informing them of the errors and asking them to correct the papers. They have not been successful. That is not my understanding of a good peer-review system. 

%If you have suggestions or would like to help, please contact me.

%If you wonder how so many erroneous papers could be published. My explanation is that humans are biased. We expect that entanglement to provide a fundamental advantage in a variety of contexts, and therefore expend little effort in checking whether this is actually the case.
\end{document}